\documentclass[letterpaper,12pt]{article}

\newcommand{\titulo}{Implementación de un relator de pantalla para sitios web basado en inteligencia artificial como herramienta de accesibilidad para personas con discapacidad visual}
\newcommand{\ciudad}{Ciudad} % e.g. Valparaíso
% TODO: Consultar el formato de los nombres:
\newcommand{\nombrealumno}{Ignacio Benjamín Riquelme Vidal}
\newcommand{\nombreprofesor}{Pedro Felipe Toledo Correa}
\newcommand{\nombrecorreferente}{Jose Luis Martí Lara}
% Mes y año del examen
\newcommand{\mesexamen}{Julio}
\newcommand{\anioexamen}{2026}
% Dedicatoria y agradecimientos
\newcommand{\dedicatoria}{
Dedico este trabajo a quienes me acompañaron en mi proceso como estudiante de principio a fin, y quienes no pudieron llegar al final:

Mi familia, por permitirme dedicarme a mis estudios sin preocupaciones externas.

Mi perrita q.e.p.d., quien desde pequeño pronostiqué me acompañaria hasta el final de mi carrera, pero por injusticias de la vida no pudo ser.

Mi pareja, por darme su apoyo incondicional en los momentos más difíciles.

Mis amigos, por traer risas cuando la vida me obligaba a tomarla en serio.

Espero tener personas —y animalitos— así de increíbles a mi lado en los próximos desafíos que la vida me depare.
}
\newcommand{\agradecimientos}{

Deseo agradecer a todos quienes aportaron en la realización de este trabajo:

Al profesor Pedro Toledo, por tomar mi idea y guiarme en su desarrollo semana a semana desde el primer día.

Al profesor Jose Luis Martí Lara, por prestar su apoyo como correferente.

Al profesor Ricardo Ñanculef, por entregar su apoyo intelectual y recursos técnicos para ejecutar el proyecto.

A José Herrera, por darme acceso a las máquinas del laboratorio de computación.

A William López, Pedro Correa, FUNDALURP y todos quienes respondieron la encuesta, quienes me ayudaron a entender la problemática de accesibilidad web desde la perspectiva de las personas ciegas.

A Luis Ortiz, por ser el primer usuario de prueba y entregar su confianza desde el primer momento.

Al doctor Alex Soto, por su apoyo para comprender la discapacidad visual desde la perspectiva de la medicina.

Sin cualquiera de estas personas, este trabajo no habría sido posible.
}
\newcommand{\resumen}{

Este trabajo presenta el desarrollo del ``Relator de pantalla'', una herramienta de accesibilidad para páginas web basada en inteligencia artificial generativa para personas con discapacidad visual. A diferencia de un lector de pantalla, que solo puede comunicar contenido textual, esta herramienta permite acceder mediante preguntas en lenguaje natural a información visual de un sitio web —imágenes, gráficos y otros elementos que un lector de pantalla no puede interpretar—, funcionando junto a él en lugar de reemplazarlo.

El desarrollo partió de una encuesta a personas en situación de discapacidad visual para validar empíricamente la necesidad, continuó con un proceso de evaluación y filtrado de los modelos de lenguaje-visión más adecuados para la tarea y finalmente el proceso de construcción de un prototipo funcional. La solución se validó mediante un conjunto de métricas técnicas y mediante entrevistas con personas ciegas en tareas reales sobre sitios web, confirmando que un asistente de este tipo puede complementar a los lectores de pantalla existentes mediante la interpretación del contenido visual, en el tipo de contenido que estos no pueden comunicar.
}
\newcommand{\resumeningles}{

This work presents the development of ``Relator de pantalla'', a web accessibility tool based on generative artificial intelligence for people with visual impairment. Unlike a screen reader, which can only communicate textual content, this tool allows users to access visual information on a website through natural language questions —images, graphics, and other elements a screen reader cannot interpret—, working alongside it instead of replacing it.

The development began with a survey of people with visual impairment to empirically validate this need, continued with a process of evaluating and filtering the vision-language models best suited for the task, and finally the construction of a functional prototype. The solution was validated through a set of technical metrics and through interviews with blind people performing real tasks on websites, confirming that an assistant of this kind can complement existing screen readers by interpreting visual content, precisely the kind of content these cannot communicate.
}
\newcommand{\palabrasclave}{

accesibilidad web; discapacidad visual; modelos de lenguaje visual (VLM); lectores de pantalla; \textit{LLM-as-a-judge}.
}
\newcommand{\palabrasclaveingles}{

web accessibility; visual impairment; vision-language models (VLM); screen readers; \textit{LLM-as-a-judge}.
}
% @@@@@@@@@@@@@@@@@@@@@@@@@@@@@@@@@@@@@@@@@@@@@@@@@@@@@@@@@@@@>

% Paquete para colores de texto (\textcolor)
\usepackage{xcolor}
% Rosado para texto guía original del template; rojo queda para anotaciones y pendientes propios
\definecolor{pink}{RGB}{255,105,180}

% Paquete para importar imágenes
\usepackage{graphicx}
% Directorio de las imágenes
\graphicspath{ {figures/} }

% Paquete para gráficos nativos de LaTeX (barras horizontales de la encuesta)
\usepackage{pgfplots}
\pgfplotsset{compat=1.18}
\usetikzlibrary{positioning,fit,arrows.meta}

% Idioma y fuentes
\usepackage[spanish,es-tabla]{babel}
\usepackage[T1]{fontenc}

\usepackage{fontspec}
% Los siguientes comandos fueron sugeridos por @anibalbastiass (ver issue#5)
% para contar con Carlito en cursiva y negrita.
\setmainfont{Carlito}[BoldFont={* Bold}]
\setmainfont{Carlito}[ItalicFont={* Italic}]

% Paquete para definir cualquier tamaño de font
\usepackage{anyfontsize}

% Settear font
\setmainfont{Carlito}

% Tamaño de la página y márgenes
\usepackage[letterpaper,top=2.5cm,bottom=3cm,left=3cm,right=3cm,marginparwidth=1.75cm]{geometry}

% Determinar interlineado:
\renewcommand{\baselinestretch}{1.0}

% Eliminar sangrías:
\setlength{\parindent}{0cm}

% Paquete para definir los formatos de los títulos
\usepackage[explicit]{titlesec}

\titleformat{name=\section}[block]{\fontsize{16}{24}\selectfont\bfseries}{}{0pt}{#1}
\titleformat{name=\section,numberless}[block]{\fontsize{16}{24}\selectfont\bfseries}{}{0pt}{#1}
\titlespacing*{name=\section}{0pt}{0pt}{0.5cm}
\titlespacing*{name=\section,numberless}{0pt}{0pt}{0.5cm}

% Separación entre parrafos
\setlength{\parskip}{0.4cm}

% Paquetes de utilidad general
\usepackage{amsmath}
\usepackage{graphicx}
\usepackage{float}
\usepackage{longtable}
\usepackage[colorlinks=true, allcolors=blue]{hyperref}

% Formato de las tablas de contenido
\usepackage{tocbasic}

% Para obtener el número de la última página
\usepackage{lastpage}

% Header y footer
\usepackage{fancyhdr}
\setlength{\headheight}{22.03pt}
\fancypagestyle{portada}{
    \lhead{}
    \chead{}
    \rhead{}
    \lfoot{}
    \cfoot{\fontsize{10}{12}\selectfont \thepage}
    \rfoot{}
    \renewcommand{\headrulewidth}{0pt}
}
\fancypagestyle{intermedio}{
    \lhead{}
    \chead{\fontsize{10}{12}\selectfont\MakeUppercase{\titulo}}
    \rhead{}
    \lfoot{}
    \cfoot{\fontsize{10}{12}\selectfont Página \textbf{\thepage}\ de \textbf{\pageref{LastPage}}}
    \rfoot{}
    \renewcommand{\headrulewidth}{1pt}
}

% Comandos para secciones
\newcommand{\secnumbersection}[1]{
\addtocounter{section}{1}
\phantomsection
\section*{CAPÍTULO \thesection \texorpdfstring{\\}\ #1}
\addcontentsline{toc}{section}{CAPÍTULO \thesection : #1}
\setcounter{subsection}{0}
}
\newcommand{\secnumberlesssection}[1]{
\section*{#1}
\phantomsection
\addcontentsline{toc}{section}{#1}
\setcounter{subsection}{0}
}

% Nombres
\addto\captionsspanish{\renewcommand{\contentsname}{ÍNDICE DE CONTENIDOS}}
\addto\captionsspanish{\renewcommand{\listfigurename}{ÍNDICE DE FIGURAS}}
\addto\captionsspanish{\renewcommand{\listtablename}{ÍNDICE DE TABLAS}}
\makeatletter
\renewenvironment{thebibliography}[1]
     {\secnumberlesssection{REFERENCIAS BIBLIOGRÁFICAS}
      \@mkboth{\MakeUppercase\bibname}{\MakeUppercase\bibname}%
      \list{\@biblabel{\@arabic\c@enumiv}}%
           {\settowidth\labelwidth{\@biblabel{#1}}%
            \leftmargin\labelwidth
            \advance\leftmargin\labelsep
            \@openbib@code
            \usecounter{enumiv}%
            \let\p@enumiv\@empty
            \renewcommand\theenumiv{\@arabic\c@enumiv}}%
      \sloppy
      \clubpenalty4000
      \@clubpenalty \clubpenalty
      \widowpenalty4000%
      \sfcode`\.\@m}
     {\def\@noitemerr
       {\@latex@warning{Empty `thebibliography' environment}}%
      \endlist}
\makeatother

% Personalizar Tabla de Contenidos

\usepackage{tocloft}
\renewcommand{\cftsecfont}{\fontsize{12}{14}\selectfont\fontspec{Carlito}}
\renewcommand{\cftsubsecfont}{\fontsize{12}{14}\selectfont\fontspec{Carlito}}
\renewcommand{\cftsubsubsecfont}{\fontsize{12}{14}\selectfont\fontspec{Carlito}}

\renewcommand\cftfigfont{\fontsize{12}{14}\selectfont\fontspec{Carlito}}

% Links sin color
\usepackage{hyperref}
\hypersetup{colorlinks = false}

% Comando para secciónes sin enumeración
% (sugerido por @anibalbastiass https://github.com/autopawn/tex-thesis-template/issues/5#issuecomment-916106128)
\newcommand{\secnumberlesssubsection}[1]{
\subsection*{#1}
\phantomsection
\addcontentsline{toc}{subsection}{#1}
\setcounter{subsection}{0}
}
% Forma de uso:
% \secnumberlesssubsection{"Sub seccion sin enumeración"}

% @@@@@@@@@@@@@@@@@@@@@@@@@@@@@@@@@@@@@@@@@@@@@@@@@@@@@@@@@@@@>
\begin{document}
\sloppy % Para evitar que referencias se escapen de los márgenes.

\pagestyle{portada}
\pagenumbering{roman}
\input{000_portadas}

\newpage
\secnumberlesssection{GLOSARIO}

{\setlength{\parskip}{0cm}

API: Interfaz de Programación de Aplicaciones (\textit{Application Programming Interface}).

ARIA: Aplicaciones Enriquecidas de Internet Accesibles (\textit{Accessible Rich Internet Applications}).

ASR: Reconocimiento Automático del Habla (\textit{Automatic Speech Recognition}).

AWS: Amazon Web Services.

CPU: Unidad Central de Procesamiento (\textit{Central Processing Unit}).

CSS: Hojas de Estilo en Cascada (\textit{Cascading Style Sheets}).

DI: Departamento de Informática.

DMAE: Degeneración Macular Asociada a la Edad.

DOM: Modelo de Objetos del Documento (\textit{Document Object Model}).

FUNDALURP: Fundación de Lucha contra la Retinitis Pigmentosa.

GB: Gigabyte.

GPU: Unidad de Procesamiento Gráfico (\textit{Graphics Processing Unit}).

HTML: Lenguaje de Marcado de Hipertexto (\textit{HyperText Markup Language}).

HTTP: Protocolo de Transferencia de Hipertexto (\textit{HyperText Transfer Protocol}).

IA: Inteligencia Artificial.

INLAB: Unidad de Inclusión Laboral (FUNDALURP).

JAWS: \textit{Job Access With Speech}.

JSON: Notación de Objetos de JavaScript (\textit{JavaScript Object Notation}).

LLM: Modelo de Lenguaje de Gran Escala (\textit{Large Language Model}).

MAE: Error Absoluto Medio (\textit{Mean Absolute Error}).

NVDA: Acceso No Visual al Escritorio (\textit{NonVisual Desktop Access}).

PS: Puntaje de Satisfacción.

RAM: Memoria de Acceso Aleatorio (\textit{Random Access Memory}).

SENADIS: Servicio Nacional de la Discapacidad.

SII: Servicio de Impuestos Internos.

SSH: Shell Segura (\textit{Secure Shell}).

SSML: Lenguaje de Marcado de Síntesis de Voz (\textit{Speech Synthesis Markup Language}).

STT: Voz a Texto (\textit{Speech to Text}), equivalente a ASR en este trabajo.

TTS: Texto a Voz (\textit{Text to Speech}), o síntesis de voz.

UC: Pontificia Universidad Católica de Chile.

UCH: Universidad de Chile.

USM: Universidad Técnica Federico Santa María (ver también UTFSM).

UTFSM: Universidad Técnica Federico Santa María.

VLM: Modelo de Lenguaje Visual (\textit{Vision-Language Model}).

VRAM: Memoria de Video (\textit{Video RAM}).

W3C: Consorcio de la World Wide Web (\textit{World Wide Web Consortium}).

WAI-ARIA: estándar del W3C para declarar la semántica de accesibilidad de elementos web (\textit{Web Accessibility Initiative -- Accessible Rich Internet Applications}).

WCAG: Pautas de Accesibilidad para el Contenido Web (\textit{Web Content Accessibility Guidelines}).
}


%Índice de contenidos:
\newpage
\thispagestyle{portada}
\tableofcontents

%Índice de figuras:
\newpage
\thispagestyle{portada}
\phantomsection
\addcontentsline{toc}{section}{ÍNDICE DE FIGURAS}
\listoffigures
\phantomsection
\addcontentsline{toc}{section}{ÍNDICE DE TABLAS}
\listoftables

\newpage
\pagestyle{intermedio}
\pagenumbering{arabic}
\secnumberlesssection{INTRODUCCIÓN}

Las personas con discapacidad visual enfrentan una dificultad particular al relacionarse con la tecnología: buena parte de su uso está fuertemente ligado a la capacidad de ver e interactuar visualmente con lo que se muestra en una pantalla. Históricamente, esta barrera se ha resuelto de distintas formas. El braille permitió traducir el texto escrito a un formato accesible mucho antes de la era digital, junto a otras soluciones como las cintas grabadas de libros o el apoyo directo de otras personas.

En el contexto digital, la forma principal en que se ha vinculado lo visual con la voz ha sido la lectura de texto: un lector de pantalla puede convertir en voz el contenido textual de una interfaz, pero solo puede apoyarse en aquello que está escrito, y no en lo que es estrictamente visual.

La capacidad de los modelos de inteligencia artificial multimodal para interpretar imágenes —determinar qué contienen, qué no contienen, e incluso razonar sobre ellas— abre la posibilidad de extender esa misma tecnología a un nuevo propósito: vincular el mundo de lo escrito con el mundo de lo no escrito. Información de un sitio web que hasta ahora era exclusivamente visual —una fotografía, un gráfico, el color o el estilo de una página— puede extraerse y verbalizarse en texto, volviéndola legible mediante las mismas tecnologías que ya hacen accesible lo escrito.

En este contexto se plantea la idea de un ``Relator de pantalla'': usar la inteligencia artificial para gestionar un relato de lo que se ve en una pantalla, asociado a las necesidades puntuales de cada persona, de modo que sea equivalente a contar con alguien al lado describiendo lo que está viendo.

Los lectores de pantalla, sin embargo, no están diseñados para interpretar este tipo de contenido, lo que abre una brecha de accesibilidad que una capa intermediaria basada en inteligencia artificial puede ayudar a cerrar —como complemento del lector de pantalla que la persona ya utiliza, no como su reemplazo—. Esta memoria caracteriza primero esa problemática, partiendo de su dimensión social y acotándola después a sus aspectos técnicos, para luego explicar cómo la inteligencia artificial ofrece una herramienta capaz de atenderla y cómo se integran específicamente los modelos de lenguaje visual para generar esa asistencia.

Para ello, el trabajo sigue una hoja de ruta que parte de una encuesta propia aplicada a personas ciegas, que motiva y valida empíricamente la problemática antes de proponer una solución; continúa con el desarrollo de un prototipo funcional y con un proceso de selección del modelo de inteligencia artificial más adecuado para la tarea; y concluye con la validación de la solución mediante métricas técnicas y pruebas con usuarios reales. El resto del documento se organiza siguiendo ese mismo recorrido: el Capítulo 1 profundiza en la problemática social y técnica que motiva este trabajo; el Capítulo 2 revisa los conceptos necesarios para comprenderlo; el Capítulo 3 describe la propuesta de solución desarrollada; el Capítulo 4 presenta su validación; y el Capítulo 5 cierra con las conclusiones y el trabajo futuro.


\newpage
\secnumbersection{DEFINICIÓN DEL PROBLEMA}

\textcolor{pink}{Se debe definir el problema, es importante no confundir definir el problema con describir la solución. Por ejemplo: ``diseñar una arquitectura e implementar una plataforma ...'' es una solución, no un problema.}

\textcolor{pink}{Algunos elementos que podrían ir en este capítulo son (no es necesario que vayan todos):
\begin{itemize}
    \item Breve descripción del contexto donde se realizará la memoria (organización, línea dentro de la Informática en la que se basa, etc.)
    \item ¿Qué y cómo se realiza actualmente la situación que mejorarás con tu memoria?
    \item ¿Qué actores o usuarios están involucrados?
    \item ¿Qué dificultades tienen esos actores actualmente? ¿cuántos son? (ideal si se pueden poner estadísticas para así saber si existe un mercado razonable para la solución que propondrás en tu memoria, en el fondo saber cuántas personas u organizaciones tienen el mismo problema que estás definiendo)
    \item ¿Qué podría pasar si en el corto o mediano plazo no se solucionan esas dificultades (¿es decir, si no se hiciera tu memoria, qué pasaría?; en el fondo justificar por qué conviene hacer tu memoria, ¿cuál es la motivación o interés de hacerla?).
    \item ¿Qué competencia existe actualmente? (a lo mejor ya existe una solución al problema, pero por qué no sirve, o por qué tu solución sería mejor, también se puede enfocar a si este problema existe en otras realidades y cómo ha sido solucionado allí).
    \item Precisar los objetivos y alcances de la memoria (o solución al problema).
\end{itemize}}

\textcolor{pink}{En este capítulo, de ser necesario puede usar referencias bibliográficas (velar porque sean recientes), una cita de ejemplo \cite{schwab2002cure} y otras más \cite{georget1994study,beaumont1990patient}.}

\textcolor{pink}{Recuerde poner notas al pie de página que sean explicativas \footnote{Este es un ejemplo de una nota al pie de página. Puede indicar alguna URL, definiciones, aclarar alguna información pertinente del texto, citar algunas referencias, etc..}.}

\subsection{PROBLEMÁTICA SOCIAL}

Las personas con discapacidad visual enfrentan barreras persistentes al interactuar con entornos digitales. La Organización Mundial de la Salud reporta que al menos 2.200 millones de personas en el mundo tienen algún grado de deficiencia visual \cite{who2026blindness}, pero esa cifra agrupa mayoritariamente casos leves y corregibles con lentes. Al filtrar por los niveles de discapacidad relevantes para esta memoria —baja visión, discapacidad visual severa y ceguera total—, las estimaciones del Global Burden of Disease para el año 2020 son de 43,3 millones de personas ciegas y 295 millones con discapacidad visual moderada a severa, es decir, un orden de 338 millones de personas a nivel mundial \cite{bourne2021trends}.

Sin embargo, la magnitud de esta cifra no es lo relevante, sino la pertinencia ética de integrar a un grupo de usuarios históricamente segregado del entorno tecnológico general, en la medida en que se ve en la obligación de depender de herramientas diseñadas exclusivamente ``para ciegos'' en lugar de acceder al mismo contenido que el resto de los usuarios. Evitar segregar o estigmatizar a un grupo de usuarios mediante soluciones separadas ha sido, de hecho, uno de los principios formales del diseño universal —el principio de uso equitativo— en los últimos 30 años \cite{centerforuniversaldesign1997}.

La consecuencia de no resolver estas barreras es la exclusión en las dimensiones social, laboral y educativa. La imposibilidad de acceder de forma autónoma a contenido en línea limita:

\begin{itemize}
    \item \textbf{Social:} la participación social cotidiana.
    \item \textbf{Laboral:} el desempeño y las oportunidades laborales.
    \item \textbf{Educativo:} el acceso a instancias educativas, que involucran a un grupo más acotado.
\end{itemize}

A modo de ejemplo, hoy no es posible enterarse por cuenta propia de qué muestra una fotografía que comparte un familiar o un amigo, ni completar sin ayuda una compra o un trámite en línea cuando de por medio se exige una verificación visual \cite{webaim2024screenreadersurvey10}. Tampoco es posible darse cuenta a tiempo de un aviso que aparece de forma inesperada en la pantalla, lo que puede significar perder una oportunidad o un plazo importante. En el trabajo, muchas veces es necesario pedirle a un compañero que describa una cifra o una imagen antes de poder tomar una decisión, y postular a un empleo puede volverse una barrera en sí misma cuando el proceso exige una verificación que solo puede completarse viendo la pantalla. Algo similar ocurre en lo educativo, donde acceder a un material de estudio compartido como imagen deja a la persona a la espera de que alguien más se lo describa antes de poder seguir aprendiendo al mismo ritmo que el resto.

\textcolor{red}{fuente pendiente: confirmar y referenciar estudio(s) específicos sobre el impacto práctico de estas barreras en el ámbito laboral y educativo.}

\subsubsection{Búsqueda de antecedentes locales}

El punto de partida se basó en que esta problemática social afecta de manera particular a las personas ciegas al navegar la web. Sin embargo, al revisar el estado del arte no se encontraron estudios específicos que caracterizaran con suficiente detalle esa necesidad en este grupo, sobre todo porque ya existen soluciones tecnológicas que resuelven parte de esto, como por ejemplo los lectores de pantalla \footnote{Programas que traducen el contenido de una página web a voz sintetizada o a un dispositivo braille; se profundiza en su funcionamiento en el Capítulo 2.}. Por ello, la problemática social se acotó hacia el contexto específico de personas con discapacidad visual que usan lector de pantalla pero aún presentan dificultades para utilizar la plenitud de la información que ofrece un sitio web al momento completar una tarea, y se diseñó y aplicó una encuesta propia con el fin de validar empíricamente la existencia y la naturaleza de la necesidad antes de proponer una solución.

\subsubsection{La encuesta (trabajo realizado)}

Se elaboró y distribuyó el instrumento ``Encuesta sobre tecnologías para personas ciegas'' mediante Google Forms, enviado por correo electrónico a miembros de la Fundación FUNDALURP (fundación de perros guía para la rehabilitación de personas en situación de discapacidad visual). La encuesta se aplicó entre febrero y marzo de 2026 y se obtuvieron 9 respondientes. El detalle completo del instrumento —su estructura, tipos de pregunta y cómo se relacionan con la problemática— se presenta en el Anexo \ref{anexo:encuesta}.

El perfil de los respondientes muestra como condición dominante la retinitis pigmentaria, con niveles de discapacidad que van desde severo hasta ceguera total, y un manejo computacional autoreportado mayoritariamente entre básico e intermedio.

Respecto de las dificultades reportadas al navegar la web con lectores de pantalla, las más frecuentes, en orden, fueron: elementos estrictamente visuales que no pueden interpretarse (imágenes, gráficos); pop-ups o notificaciones que no se leen; campos de escritura sin etiqueta; la obligación de usar el ratón para ciertas interacciones; y los captchas visuales. \textcolor{red}{pendiente: agregar conteo o porcentaje exacto de respondientes que reportó cada dificultad.}

En cuanto a las funcionalidades más demandadas, se destacan tres, en orden de frecuencia: describir e interpretar imágenes; la navegación por voz hacia un objetivo concreto; y la lectura de notificaciones o etiquetas dinámicas. \textcolor{red}{pendiente: agregar el conteo o porcentaje exacto de respondientes que solicitó cada funcionalidad.}

\textcolor{red}{gráficos: distribución del perfil de respondientes (condición, nivel de discapacidad, nivel computacional) y frecuencia de dificultades y demandas reportadas.}

\subsection{PROBLEMÁTICA TÉCNICA}

Actualmente, las personas con discapacidad visual navegan la web apoyándose en lectores de pantalla, programas como JAWS, NVDA o VoiceOver que traducen el contenido de una página a voz sintetizada o a un dispositivo braille. El DOM (Document Object Model) es la estructura de nodos que el navegador construye a partir del código HTML de la página; a partir de ella, el navegador deriva un árbol de accesibilidad —una representación paralela que expone únicamente los roles, nombres accesibles, estados y relaciones relevantes para una tecnología asistiva— y lo pone a disposición a través de la API de accesibilidad del sistema operativo \cite{mdn2024accessibilitytree}. Es esa API la que el lector de pantalla consulta para construir lo que finalmente comunica al usuario. También existe la posibilidad de que el usuario navegue por una página utilizando únicamente el teclado, lo cual es posible gracias al atributo \texttt{tabindex} y al orden del DOM. El recorrido secuencial mediante la tecla Tab lo controla el propio navegador para la navegación por teclado, en el orden del DOM o según el atributo \texttt{tabindex}, y no es exclusivo de los lectores de pantalla ni de los usuarios con discapacidad visual.

Buena parte de la información que un lector de pantalla puede comunicar depende de las etiquetas ARIA (Accessible Rich Internet Applications): un conjunto de atributos que se agregan a los elementos del DOM para declarar roles, estados y relaciones que no tienen una etiqueta HTML nativa que los describa \cite{mdn2024aria}. Por ejemplo, una notificación emergente puede marcarse como una región viva (\texttt{aria-live}) para que el lector la anuncie apenas aparece. Cuando un sitio no declara estas etiquetas, el lector de pantalla no tiene forma de identificar ese contenido dinámico, aunque técnicamente sí forme parte del DOM.

El límite de este enfoque está, entonces, en que depende por completo de que la página declare la semántica correcta —HTML nativo y ARIA— para cada elemento. Los lectores de pantalla no interpretan información estrictamente visual: no describen imágenes, gráficos ni captchas, y solo leen de forma confiable las notificaciones o pop-ups que fueron etiquetados correctamente por quien construyó el sitio. Tampoco comprenden la estructura espacial de la interfaz (qué elemento está dónde y con qué propósito) más allá del orden en que fueron declarados. El usuario recibe así una representación incompleta de la página, que depende de decisiones de accesibilidad ajenas a su control.

Respecto de la competencia y soluciones existentes, las alternativas actuales resultan insuficientes. Los lectores de pantalla, por diseño, no cubren la interpretación visual; y las soluciones basadas en IA para descripción de imágenes no están integradas al flujo de navegación ni fueron pensadas específicamente para personas ciegas, por lo que no resuelven el problema de forma completa.

\textcolor{red}{tabla: comparación de soluciones existentes (lectores de pantalla, herramientas de descripción de imágenes por IA etc) frente a las capacidades requeridas, para justificar la brecha que cubre la propuesta.}

\textcolor{red}{fuente pendiente: referencias recientes sobre lectores de pantalla, accesibilidad web (p. ej. WCAG) y soluciones de IA para descripción de imágenes.}

\subsection{OBJETIVOS Y ALCANCES}

Considerando la relevancia de que las personas con discapacidad visual puedan acceder de forma equitativa a las tecnologías de la información, se plantean los siguientes objetivos para el desarrollo del asistente de accesibilidad web que en este trabajo se denomina Relator.

\textbf{Objetivo general.} Desarrollar el Relator, un sistema basado en inteligencia artificial capaz de analizar e interpretar la información visual de sitios web a partir de capturas de pantalla, utilizando modelos de lenguaje visual (VLM), que se integre y coexista con el lector de pantalla del usuario, con el fin de mejorar la accesibilidad de las personas con discapacidad visual.

\textbf{Objetivos específicos.}

\begin{itemize}
    \item[(a)] Implementar un prototipo funcional del asistente (Relator) capaz de interpretar información visual y entregar descripciones auditivas.
    \item[(b)] Definir instrucciones de navegación que orienten al modelo en la interpretación de los elementos web.
    \item[(c)] Empaquetar e implementar la solución en un entorno accesible y adaptable.
    \item[(d)] Validar la solución mediante pruebas con usuarios reales con discapacidad visual.
\end{itemize}

\textbf{Alcances.} \\Se espera que la aplicación funcione sobre navegadores basados en Chromium (Google Chrome, Microsoft Edge, Opera, entre otros) cubriendo la mayoría de los casos de uso. \textcolor{red}{(Incluir cifras de uso de navegadores)} El idioma de interacción contemplado es el español, aunque se espera que el modelo VLM seleccionado soporte de forma nativa otros idiomas \textcolor{red}{(Esto se confirmará en Capitulo 3. Propuesta de Solución)}.

\newpage
\secnumbersection{MARCO CONCEPTUAL}

\subsection{LECTORES DE PANTALLA Y EL DOM}

Un lector de pantalla es un software de tecnología asistiva que convierte el contenido de una interfaz digital en voz sintetizada o en salida braille, permitiendo al usuario acceder a esa información sin necesidad de verla \cite{w3ctoolstechniques}. La salida braille se entrega mediante una línea braille refrescable: un dispositivo de hardware, conectado al computador, con celdas que forman y deshacen caracteres en relieve mediante pines accionados electromecánicamente, y que va mostrando al tacto y en tiempo real el contenido que el lector de pantalla comunica.

Los lectores de pantalla no son una categoría homogénea ni universal: existen embebidos en el propio sistema operativo, como Windows Narrator o VoiceOver en macOS/iOS, y otros que se instalan por separado, como JAWS o NVDA; dentro de estos últimos hay opciones de código abierto y sin costo (NVDA) y opciones pagas orientadas al uso profesional (JAWS); y hay tanto lectores de escritorio como lectores para dispositivos móviles, como TalkBack en Android. Qué lector de pantalla resulta más adecuado depende, entonces, del sistema operativo, el dispositivo y el presupuesto disponible, y no existe uno que cubra todos los contextos de uso por igual.

\begin{table}[h]
    \centering
    \caption{\label{table:lectorespantalla} Clasificación de los lectores de pantalla más usados.} Fuente: Elaboración propia, a partir de la Encuesta de Usuarios de Lectores de Pantalla \#10 de WebAIM \cite{webaim2024screenreadersurvey10}.
    \begin{tabular}{|p{2.4cm}|p{2.3cm}|p{2.8cm}|p{2.3cm}|p{4.2cm}|}
        \hline
        \textbf{Lector} & \textbf{Tipo} & \textbf{Licencia} & \textbf{Plataforma} & \textbf{Uso reportado} \\
        \hline
        JAWS & Instalado & Pago & Escritorio & 40,5\% como lector principal \\
        \hline
        NVDA & Instalado & Código abierto, sin costo & Escritorio & 37,7\% como lector principal \\
        \hline
        VoiceOver & Embebido (macOS/iOS) & Incluido con el sistema operativo & Escritorio y móvil & 9,7\% como lector principal de escritorio; 70,6\% de uso reportado en móviles \\
        \hline
        Narrator & Embebido (Windows) & Incluido con el sistema operativo & Escritorio & 0,7\% como lector principal \\
        \hline
        TalkBack & Embebido (Android) & Incluido con el sistema operativo & Móvil & 34,7\% de uso reportado en móviles \\
        \hline
    \end{tabular}
\end{table}

Las cifras de escritorio corresponden a la pregunta por el lector de pantalla \textbf{principal}, mientras que las cifras de uso en dispositivos móviles corresponden a una pregunta distinta por los lectores que el encuestado usa \textbf{con frecuencia}, sin exclusividad entre ellos \cite{webaim2024screenreadersurvey10}; por eso no son directamente comparables entre sí ni suman 100\%.

El DOM (\textit{Document Object Model}) es la representación estructurada del documento web: una jerarquía de nodos (elementos, atributos, texto) que el navegador construye a partir del código HTML de la página en su parte estática inicial, y sobre la que se puede consultar y operar mediante programación \cite{whatwgdom}. Esa segunda dimensión, programática, es la que aporta JavaScript: al ejecutarse, puede agregar, quitar o modificar nodos del DOM ya construido, de modo que lo que el navegador termina exponiendo no depende solo del HTML servido inicialmente, sino también de los cambios que JavaScript introduce después. El navegador deriva del DOM resultante un árbol de accesibilidad, una representación paralela y más acotada que expone únicamente los roles, nombres, estados y relaciones relevantes para una tecnología asistiva, a través de la API de accesibilidad del sistema operativo (por ejemplo, UI Automation en Windows o AXAPI en macOS). Es esa API la que el lector de pantalla consulta para construir lo que finalmente comunica al usuario.

La capacidad del navegador de construir un árbol de accesibilidad completo depende de que la página declare correctamente su semántica, algo que el HTML sin marcadores adicionales no siempre garantiza. El W3C define para ello dos estándares complementarios: las Pautas de Accesibilidad para el Contenido Web (WCAG), que establecen los criterios que debe cumplir un sitio para considerarse accesible \cite{wcag22}, y WAI-ARIA, un conjunto de atributos (roles, estados y propiedades \texttt{aria-*}) que permite declarar la semántica de elementos que no la expresan por sí solos —por ejemplo, marcar una notificación emergente como una región viva (\texttt{aria-live}) para que se anuncie apenas aparece— \cite{waiaria2023}. El atributo \texttt{alt}, por su parte, es el mecanismo nativo para asociar un texto alternativo a una imagen —y existe desde antes que los lectores de pantalla, no exclusivamente para ellos—; sin él, una imagen no tiene ninguna representación textual que el lector de pantalla pueda comunicar \cite{w3caltimages}.

El lector de pantalla permite al usuario recorrer \textbf{secuencialmente} el contenido de la página mediante la tecla Tab, uno tras otro: no solo los elementos interactuables como campos de texto o enlaces, sino también contenido no interactuable como párrafos de texto. Ese orden de recorrido no es estrictamente el orden en que los elementos fueron declarados en el HTML: el atributo \texttt{tabindex} permite alterarlo, adelantando o postergando elementos respecto de su posición original en el documento. Además de ese recorrido, el lector de pantalla ofrece ayudas de navegación que no existen sin él: una de las más usadas es el listado de enlaces, que presenta todos los enlaces de la página en una lista y permite saltar directamente al que se necesite.

La eficacia de este recorrido secuencial depende de la familiaridad previa del usuario con la estructura de la página: sobre un sitio ya explorado, permite ubicar rápidamente los elementos buscados; frente a un sitio nuevo, requiere recorrerlo y estudiarlo varias veces antes de poder anticipar su estructura.

\subsection{EXTENSIONES DE NAVEGADOR}

Una extensión es un módulo de software que un usuario instala sobre su navegador para modificar o ampliar su comportamiento —por ejemplo, para bloquear anuncios, gestionar contraseñas o capturar información de una página—, sin ser una aplicación independiente ni una función nativa del propio navegador \cite{mdnwebextensions}.

Google Chrome fue el primer navegador en incorporar un sistema de extensiones basado enteramente en HTML, CSS y JavaScript, habilitado en su canal de desarrollo en 2009 y lanzado de forma pública junto a Chrome 4.0 en enero de 2010 \cite{chromeextensionshistory}. A partir de esa base, Mozilla adoptó en 2015 una API compatible para su navegador Mozilla Firefox, bautizada WebExtensions, que buscaba convertirse en un estándar común entre navegadores; en 2021 se llegó a formar un grupo comunitario del W3C específicamente para intentar formalizar esa base común entre distintos fabricantes \cite{w3cwebextensionscg2021}, pero ese intento de estandarización no llegó a concretarse. Pese a ello, la popularidad de Chrome llevó a que el resto de los navegadores basados en su motor Chromium (Edge, Opera, Brave, entre otros) convergieran hacia una API de extensiones muy similar \cite{mdnwebextensions}.

Hoy existen, en términos generales, dos grandes familias de navegador incompatibles entre sí: los basados en Chromium, que agrupan a la gran mayoría, y Mozilla Firefox, basado en el motor Gecko y heredero del antiguo navegador Netscape. Aparte de estas dos familias existen otros navegadores, como los propietarios de ciertos fabricantes de dispositivos móviles. Un caso aparte es Safari que utiliza su propio motor, WebKit. Independiente de que un navegador soporte extensiones, estas solo son compatibles con el navegador para el que fueron diseñadas: el código y el empaquetado de una extensión para Chrome no son directamente instalables en Firefox ni viceversa, aún cuando ambas compartan la mayoría de las llamadas de su API.

A nivel de capacidades, una extensión se ejecuta bajo un modelo de permisos explícitos declarados en un archivo de manifiesto: solo puede acceder a las páginas, pestañas o funciones del navegador para las que fue autorizada, y se organiza mediante scripts de contenido —que se ejecutan en el contexto de la página visitada— y un script de fondo, que coordina la lógica de la extensión y sus comunicaciones con servicios externos. Ese modelo de permisos delimita tanto lo que la extensión puede automatizar dentro del navegador como lo que queda fuera de su alcance sin intervención directa del usuario o del sistema operativo.

\subsection{RECONOCIMIENTO Y SÍNTESIS DE VOZ (STT Y TTS)}

El reconocimiento automático del habla (\textit{automatic speech recognition}, ASR; también conocido como \textit{speech to text}, STT) es el proceso mediante el cual un sistema decodifica una señal de audio con voz humana y la transcribe a texto, típicamente combinando un modelo acústico —que traduce la señal a una secuencia probable de fonemas— y un modelo de lenguaje —que determina, a partir de esos fonemas, las palabras más probables— \cite{ibmasr}.

La síntesis de voz (\textit{text to speech}, TTS) es el proceso inverso: toma una entrada textual y la convierte en audio hablado. Esa entrada no siempre es texto plano: estándares como SSML (\textit{Speech Synthesis Markup Language}), del W3C, permiten anotar el texto con marcadores de énfasis, acento, velocidad o pausas, para un control más fino de cómo se pronuncia \cite{w3cssml}. A nivel de técnica, la síntesis de voz se ha resuelto históricamente de más de una forma: por concatenación de fragmentos de voz pregrabada, por síntesis paramétrica a partir de un modelo estadístico de la voz, y, más recientemente, mediante modelos neuronales generativos que producen la forma de onda directamente a partir de datos de entrenamiento \cite{tan2021neuraltts}.

Lectores de pantalla como JAWS o NVDA no necesariamente implementan su propio motor de síntesis: en general, delegan la conversión final de texto a voz a un motor TTS del sistema operativo o a una voz de terceros instalada por el usuario, y lo que le corresponde al lector de pantalla es determinar qué texto de la interfaz enviar a ese motor.

\subsection{IA MULTIMODAL: MODELOS DE LENGUAJE VISUAL (VLM)}

\subsubsection{IA generativa de lenguaje}

Los modelos de lenguaje generativos son un tipo de inteligencia artificial capaz de producir texto nuevo a partir de patrones aprendidos de grandes volúmenes de datos, en lugar de limitarse a clasificar o predecir una etiqueta sobre texto ya existente \cite{ibmgenpredictive}. Esa generación se implementa técnicamente como una predicción: un modelo autoregresivo —que genera cada nueva unidad de texto a partir de lo que él mismo generó previamente—, como los que se basan en la arquitectura Transformer —una arquitectura de red neuronal que pondera mediante mecanismos de atención qué partes de la entrada son más relevantes para cada predicción—, no produce una respuesta completa de una sola vez: predice, \textit{token} a \textit{token} —el \textit{token} es la unidad mínima de texto en que el modelo divide el lenguaje, por ejemplo una palabra o parte de ella—, cuál es la unidad de texto más probable dado todo lo que ha recibido y generado hasta ese momento, y repite ese proceso hasta completar la respuesta \cite{vaswani2017attention}. Es decir, ``generar'' texto y ``predecir la palabra siguiente'' no son dos capacidades distintas, sino una relación de mecanismo a resultado: el modelo predice de forma iterativa, \textit{token} a \textit{token}, y esa cadena de predicciones es percibida por quien lo usa como la generación de una respuesta nueva.

\subsubsection{IA multimodal}

La IA multimodal corresponde a modelos de IA capaces de procesar y/o producir información de más de un tipo de dato —texto, imágenes, audio o video— de forma conjunta, en lugar de limitarse a un único tipo de entrada o salida \cite{wu2023multimodalsurvey}. Dentro de esa familia se encuentran los modelos de lenguaje visual (VLM, \textit{vision language model}), especializados en la interacción entre texto e imágenes: reciben como entrada una combinación de imagen(es) y texto, y devuelven una salida en texto, apoyándose en una representación interna que alinea los conceptos visuales de la imagen con los conceptos del lenguaje \cite{wu2023multimodalsurvey}.

\subsubsection{Entrenamiento}

El entrenamiento es el proceso mediante el cual un modelo de IA aprende a reconocer patrones y correlaciones para optimizar su desempeño en tareas similares a las de su uso previsto \cite{ibmmodeltraining}. Ese aprendizaje se realiza sobre un \textit{dataset}: una colección de datos de ejemplo relevantes para la tarea que el modelo debe resolver —por ejemplo, pares de pregunta y respuesta, o imágenes junto a su descripción—; cuanto más se asemeje ese conjunto a los casos reales que el modelo enfrentará después, mejor será su desempeño sobre datos nuevos \cite{ibmtrainingdata}.

Un entrenamiento puede fallar en dos direcciones opuestas. Por sobreajuste (\textit{overfitting}), el modelo se ajusta tan estrechamente a los ejemplos de entrenamiento que termina memorizándolos en lugar de generalizar, logrando buenos resultados solo sobre datos ya vistos; por subajuste (\textit{underfitting}), el modelo es demasiado simple, o el entrenamiento demasiado breve, para siquiera capturar los patrones del propio conjunto de entrenamiento \cite{ibmoverfitting}. Un entrenamiento adecuado, en cambio, encuentra un punto intermedio entre ambos extremos, que le permite generalizar razonablemente bien a datos nuevos de una naturaleza similar a los ya vistos.

De ahí se desprende una limitación central de cualquier modelo de IA: no puede responder de forma confiable sobre información que su \textit{dataset} de entrenamiento no consideró, o que consideró en una cantidad insuficiente para generalizar ese patrón sin sesgo. Cuando un tipo de caso está subrepresentado en el \textit{dataset}, el modelo tiende a fallar o a generalizar incorrectamente para ese caso a partir de lo aprendido sobre los casos mayoritarios \cite{ibmdatabias}. En términos simples, un modelo entrega mejores resultados sobre aquello que aparece con frecuencia en su \textit{dataset} de entrenamiento, y resultados progresivamente peores sobre aquello que aparece con poca frecuencia o no aparece en él.

\subsubsection{Inferencia}

La inferencia es el acto de usar un modelo de IA ya entrenado para obtener una salida a partir de una entrada nueva \cite{nvidiaaiinference}. Ejecutar esa inferencia sobre un modelo de gran tamaño es, en sí mismo, un desafío computacional: un equipo de escritorio promedio no puede hacerlo, o al menos no en tiempos que permitan un uso práctico, por lo que habitualmente se recurre a hardware específico —típicamente GPU— para hacerla eficiente. Frente a esto existen dos caminos habituales, ambos con inconvenientes: los servicios de IA en línea corren sobre datacenters de gran escala y cobran por consulta; además, procesan los datos del usuario en infraestructura de terceros, sin garantías claras sobre su tratamiento.

En este contexto aparece el concepto de latencia: el tiempo de espera entre que se envía una consulta y se recibe la respuesta. Es un factor crítico para la experiencia de uso, ya que una espera prolongada puede interrumpir el flujo de la interacción del usuario, incluso cuando el contenido de la respuesta es correcto.

\subsection{EVALUACIÓN AUTOMATIZADA DE VLMs}

Existe una gran oferta de modelos VLM \textit{open source}, y distintos modelos entregan resultados diferentes para las mismas consultas, por lo que elegir el más pertinente para una necesidad determinada no es trivial. Hacerlo de forma manual —tomar un VLM, ingresarle N preguntas y evaluar cada una de esas N respuestas— no escala: para que el resultado sea confiable se necesita un N grande, y probar varios modelos vuelve la evaluación manual una tarea repetitiva e inviable. De ahí la necesidad, documentada en la literatura, de automatizar este proceso mediante técnicas de \textit{benchmarking} sistemático.

\subsubsection{Rúbrica cualitativa}

Una rúbrica es un instrumento de evaluación cualitativa que describe, para cada nivel de desempeño posible, los criterios concretos que una respuesta debe cumplir para alcanzarlo; se construye de forma manual antes de la evaluación y permite calificar de forma consistente incluso cuando la respuesta es abierta y no tiene una única forma correcta —por ejemplo, al calificar la respuesta de un estudiante en un examen de argumentación, donde una rúbrica puede distinguir niveles que van desde una postura sin fundamento hasta un argumento completo y bien sustentado. La evaluación de un VLM enfrenta el mismo problema de fondo: la respuesta a una consulta sobre una imagen tampoco tiene una única forma correcta, por lo que una rúbrica permite distinguir niveles de calidad —por ejemplo, si identifica los elementos relevantes de la imagen, si omite información importante, o si incluye contenido que no está presente en ella— en lugar de exigir una coincidencia exacta con una respuesta de referencia.

\subsubsection{\textit{LLM-as-a-judge}}

``\textit{LLM-as-a-judge}'' es un paradigma de evaluación en el que un LLM, en lugar de un evaluador humano, cumple el rol de juez: se le entrega la respuesta que se quiere evaluar junto con los criterios de evaluación —por ejemplo, una rúbrica—, y el modelo determina si esos criterios están presentes o ausentes para asignar un puntaje de forma cuantitativa \cite{zheng2023llmasjudge}.

\subsubsection{\textit{LLM-as-a-judge-evaluation}}

Para confirmar que un LLM evaluador funciona bien, en este trabajo se introduce el concepto de ``\textit{LLM-as-a-judge-evaluation}'': un humano califica de forma independiente las mismas respuestas que evalúa el LLM, y luego se contrasta cuánto se acercan los puntajes de uno y otro. Es una segunda capa de evaluación, destinada a comprobar que el criterio con el que el LLM juzga esas respuestas está bien calibrado.

Con la \textit{LLM-as-a-judge-evaluation} se cierra el ciclo: el humano evalúa las respuestas, el LLM las evalúa también, y la comparación entre ambos permite validar al LLM evaluador a partir del criterio humano.

\subsubsection{Herramientas estandarizadas de evaluación}

Respecto de los sistemas estandarizados de evaluación específica para VLMs, ya existen herramientas y métricas destinadas a esta tarea. Un ejemplo es VLMEvalKit (open-compass): un \textit{toolkit open source} de evaluación de modelos multimodales que unifica la evaluación de una gran cantidad de modelos sobre decenas de \textit{benchmarks} con configuración mínima, automatizando la preparación de datos, la inferencia y el cálculo de métricas, e incluso usando un LLM como extractor de respuestas cuando la comparación exacta falla \cite{duan2024vlmevalkit}. Sin embargo, ninguna de estas herramientas está diseñada específicamente para la lectura de sitios web en el contexto de la discapacidad visual, que es el foco de este trabajo.
\newpage
\secnumbersection{PROPUESTA DE SOLUCIÓN}

\textcolor{pink}{Se debe desarrollar la solución propuesta. Los subcapítulos por poner aquí son propios del autor. Se sugiere mencionar metodología usada. Es conveniente incorporar figuras y tablas para aclarar la solución, que deben indicar el número de la figura, su nombre y su autor o fuente (si las diseñas tú, la fuente es ``Elaboración propia''). Ver ejemplos en esta página y en la siguiente.}

\textcolor{pink}{Cabe mencionar que aquí está la esencia del trabajo en lo que se refiere al aporte creativo del memorista, es el momento de demostrar que usted es un destacado profesional que creó, diseñó y/o llevó a cabo la solución propuesta.}

\textcolor{pink}{\secnumberlesssubsection{EJEMPLO DE COMO CITAR FIGURAS E ILUSTRACIONES}}

\textcolor{pink}{Se colocó una imagen que se puede referenciar también desde el texto (Ver figura \ref{fig:malla}).}

\begin{figure}[h]
\centering
\includegraphics[width=0.8\textwidth]{malla_ingenieria_informatica}
\caption{\label{fig:malla} \textcolor{pink}{Malla Curricular Ingeniería Civil Informática.}} \textcolor{pink}{Fuente: Departamento de Informática.}
\end{figure}

\subsection{ARQUITECTURA GENERAL DEL RELATOR}

\begin{itemize}
    \item Arquitectura del Relator: sitio -> extensión (captura) -> maqueta (voz + prompt) -> albert (VLM) -> maqueta (TTS) -> audio.
\end{itemize}

\subsection{EXTENSIÓN DE CAPTURA (ss-extension)}

\begin{itemize}
    \item Extensión (ss-extension): captura automática en Chrome, viewport visible y página completa (Scroll \& Stitch); output JPG o base64.
\end{itemize}

\subsection{MAQUETA / CONTROLADOR (ss-maqueta)}

\begin{itemize}
    \item Maqueta (ss-maqueta): controlador que orquesta captura + voz + prompt, envía a albert y sintetiza la respuesta en audio.
    \begin{itemize}
        \item 3 iteraciones, cada una una mini-app funcional por sí sola:
        \item textToText: envía texto al modelo y recibe respuesta en texto.
        \item textToVoice: envía texto y recibe TTS de vuelta, integrado con lectores NVDA y JAWS.
        \item voiceToVoice: input por voz vía STT (Web Speech API de Google) y salida por TTS del lector de pantalla.
        \item La última iteración engloba a la anterior mediante comandos para alternar entre input por texto y por voz.
        \item Naturaleza: prototipo funcional para validar el pipeline completo, no el producto final.
    \end{itemize}
    \item Objetivo (b): las instrucciones de navegación se materializan en el diseño de prompts del sistema y del benchmark (ver \ref{sec:benchmark}).
\end{itemize}

\subsection{INFRAESTRUCTURA Y SERVIDOR ALBERT}

\begin{itemize}
    \item albert: servidor GPU institucional que ejecuta los las consultas a VLM; acceso vía jump host.
    \item Decisión de infraestructura: albert sobre local (GPU inviable para desarrollador ni usuario), cloud (cold start 10–20 s) y APIs (costo/privacidad).
    \item Latencia objetivo: <15 s por consulta; referencia InternVL2-2B ~1 min en T4 de Google Colab.
\end{itemize}

\subsection{BENCHMARK DE MODELOS VLM (ss-benchmark)}
\label{sec:benchmark}

\begin{itemize}
    \item Modelos candidatos: moondream2, InternVL2-2B, InternVL3-2B, Qwen3-VL-8B; compromiso entre razonamiento y hardware (necesidad de albert).
    \item Benchmark (ss-benchmark): Procedimiento para seleccionar el VLM más adecuado evaluando sobre sitios reales.
    \begin{itemize}
        \item Corpus: 10 sitios en 4 categorías (educacional, estatal, e-commerce, foros).
        \item Preguntas: 6 por sitio (3 generales + 3 específicas), 60 en total con rúbrica 0–3 definida por pregunta; 60 consultas por modelo.
        \item Rúbrica 0–3: 3 exacta, 2 incompleta, 1 parcial, 0 alucinación; rigor literal y sin sesgo de verbosidad.
        \item Pipeline: screenshots -> VLM -> JSON -> juez LLM -> métricas.
        \item Metodología: enfoque secuencial (implementado) y combinatorio (propuesto, hasta 43.200 consultas) contra sesgo de orden.
        \item Criterios de selección: precisión espacial, latencia, razonamiento semántico, resistencia a alucinaciones.
    \end{itemize}
\end{itemize}

\newpage
\secnumbersection{VALIDACIÓN DE LA SOLUCIÓN}

\textcolor{pink}{Se debe validar la solución propuesta. Esto significa probar o demostrar que la solución propuesta es válida para el entorno donde fue planteada.}

\textcolor{pink}{Tradicionalmente es una etapa crítica, pues debe comprobarse por algún medio que vuestra propuesta es básicamente válida. En el caso de un desarrollo de software es la construcción y sus pruebas; en el caso de propuestas de modelos, guías o metodologías podrían ser desde la aplicación a un caso real hasta encuestas o entrevistas con especialistas; en el caso de mejoras de procesos u optimizaciones, podría ser comparar la situación actual (previa a la memoria) con la situación final (cuando la memoria está ya implementada) en base a un conjunto cuantitativo de indicadores o criterios.}

\textcolor{pink}{\subsection{EJEMPLO DE COMO CITAR TABLAS}}

\textcolor{pink}{Se colocó una tabla que se puede referenciar también desde el texto (Ver tabla \ref{table:coloquios}).}

\begin{table}[h]
    \centering
    \caption{\label{table:coloquios} \textcolor{pink}{Coloquios del Ciclo de Charlas Informática.}} \textcolor{pink}{Fuente: Elaboración Propia.}
    \begin{tabular}{|p{7cm}|p{7cm}|}
        \hline
        \textcolor{pink}{Título Coloquio} & \textcolor{pink}{Presentador, País} \\
        \hline
        \textcolor{pink}{``Sensible, invisible, sometimes tolerant, heterogeneous, decentralized and interoperable... and we still need to assure its quality...'''} & \textcolor{pink}{Guilherme Horta Travassos, Brasil.}\\
        \hline
        \textcolor{pink}{``Dispersed Multiphase Flow Modeling: From Environmental to Industrial Applications'''} & \textcolor{pink}{Orlando Ayala, EE.UU.}\\
        \hline
        \textcolor{pink}{``Líneas de Producto Software Dinámicas para Sistemas atentos el Contexto'''} & \textcolor{pink}{Rafael Capilla, España.}\\
        \hline
        \textcolor{pink}{...} & \textcolor{pink}{...} \\
        \hline
    \end{tabular}
\end{table}

\subsection{ENFOQUE DE LA VALIDACIÓN}

A fin de validar el desarrollo, se realizaron pruebas con usuarios reales con discapacidad visual sobre el sistema ya construido. Se busca medir qué tan bien el Relator resuelve tareas reales de interpretación de imágenes en sitios web, integrando las consultas al Relator como parte del flujo de trabajo habitual del usuario con su lector de pantalla.

Dada la índole innovadora del trabajo, resulta difícil realizar una comparación cuantitativa formal —por ejemplo, contrastar los resultados de las tareas contra el lector de pantalla—, ya que los resultados esperados de uno y otro difieren demasiado entre sí. Por ello se opta por una prueba directa con el usuario, consultando su nivel de satisfacción tras completar cada tarea, llevando algo cualitativo (la sensación de satisfacción) a un valor cuantificable (un número entero).

\subsection{DISEÑO DEL PROCEDIMIENTO DE PRUEBA}

Se establece un número objetivo de entre 3 y 5 usuarios. Se exige que todos presenten discapacidad visual severa o ceguera total, que sean usuarios habituales de lectores de pantalla y que tengan un nivel de manejo computacional intermedio o avanzado.

Como entorno de pruebas se utiliza un subconjunto del corpus del benchmark (ver \ref{sec:benchmark}): al menos 2 sitios de categorías distintas y al menos 2 tareas por cada sitio. Además, se exige utilizar al menos una vez la modalidad por voz (VTV).

El detalle de los resultados por tarea se presenta en una tabla en el anexo. \textcolor{red}{REFERENCIA PENDIENTE: insertar \ref{...} a la tabla del anexo correspondiente.}

\subsection{MÉTRICA E INSTRUMENTO DE MEDICIÓN}

\subsubsection{Puntaje de satisfacción del usuario}

Como se mencionó, al terminar el flujo de trabajo de cada tarea se le pregunta al usuario su nivel de satisfacción. Este es un número entero entre 1 y 5, denominado puntaje de satisfacción (PS), donde 1 corresponde a insatisfacción total (el Relator alucinó información o su respuesta fue de nula ayuda para completar la tarea), 3 a satisfacción parcial (el Relator cumplió con lo mínimo solicitado, pero hay espacio de mejora) y 5 a satisfacción total (la respuesta fue la esperada y ayudó a completar la tarea sin problemas).

\subsubsection{Criterio de completitud}

Se definen como tareas parciales o completadas todas aquellas que obtengan un PS mayor o igual a 3. Este criterio se comenta a los usuarios al momento de consultar su puntaje, por lo que se asume una correlación directa entre el valor entregado y la completitud de la tarea.

\subsubsection{Estándar mínimo}

Se considera una sesión aprobada si, sobre el total de PS, se obtiene al menos un 55\% de tareas parciales o completadas (PS $\geq$ 3) y un promedio de al menos 55\%, equivalente a un PS promedio $\geq$ 2,75. El estándar mínimo por sesión se extrapola a los resultados globales (sumando todas las tareas de todas las sesiones y calculando sobre el total). Por lo tanto, para validar la solución propuesta se espera obtener al menos un 55\% de tareas aprobadas y un PS promedio global mayor a 2,75.

\textcolor{red}{figura: fórmula matemática para el cálculo de métricas.}


\subsection{EJECUCIÓN}

\subsubsection{Participantes}

La validación se realizó finalmente con 3 participantes, dentro del rango objetivo definido. Los tres presentan ceguera total y usan JAWS como lector de pantalla. Se diferencian entre sí en las condiciones asociadas a su discapacidad visual y en su nivel de manejo computacional (uno con nivel medio y dos con nivel avanzado).

\subsubsection{Sesiones}

Sobre estos 3 participantes se realizaron 4 sesiones en total, ya que P2 realizó dos sesiones. De ellas, 3 se hicieron en modalidad texto (TTV) y 1 en modalidad voz (VTV).

\textcolor{red}{tabla: detalle de sesiones (acá o anexo?).}

\subsection{RESULTADOS GLOBALES}

En total se evaluaron 19 tareas a lo largo de las 4 sesiones. El promedio global de satisfacción fue de 2,84 sobre 5 (aproximadamente 56\%), y 11 de las 19 tareas (57,9\%) obtuvieron un puntaje igual o superior a 3. Estos resultados cumplen ambos umbrales definidos como estándar mínimo en la sección anterior.

\textcolor{red}{tabla: resultados globales de satisfacción por tarea y sesión.}

\subsection{ANÁLISIS DE PATRONES}

A partir de la revisión de los registros de las sesiones se identificaron patrones en el desempeño del sistema. Por categoría de sitio, el mejor desempeño se observó en sitios de e-commerce y redes sociales, mientras que el peor desempeño se dio en sitios educacionales, particularmente en portales de gestión interna que requieren autenticación. Por tipo de tarea, el sistema tuvo buen desempeño en tareas descriptivas o de extracción de información --por ejemplo, ``qué es'', ``qué hay'' o ``describe''--, y un desempeño deficiente en tareas procedimentales o de navegación --por ejemplo, ``cómo hago'', ``dónde voy'' o solicitudes de pasos a seguir--.

Entre los modos de falla identificados se cuentan: un sesgo de perspectiva vidente en las respuestas del modelo, alucinaciones cuando el elemento consultado no está visible en la captura, descripciones superficiales, y fallas del mecanismo de scroll \& stitch (ver Capítulo 3, sección de extensión de captura) al capturar ventanas emergentes (no pop-ups, sino ventanas independientes usualmente para una sola función, como por ejemplo pagos online o visualizar documentos).

En la modalidad de voz se observó un comportamiento particular: se reportó una satisfacción más alta en la misma tarea ejecutada por texto incluso cuando esta no se completaba. Esto sugiere que la experiencia de interacción por voz tiene un valor intrínseco, independiente del éxito de la tarea.

\subsection{RETROALIMENTACIÓN DE LOS USUARIOS}

Los participantes aportaron las siguientes sugerencias de mejora durante las sesiones:

\begin{itemize}
    \item Contar con una versión móvil y compatible con aplicaciones, no solo con navegadores de escritorio.
    \item Que el chat mantuviera memoria de la conversación, para poder seguir preguntando sobre un mismo elemento sin repetir contexto.
    \item Poder acceder o navegar directamente a un enlace ya identificado por el sistema.
    \item Mejorar el feedback del sistema de entrada por voz, para que el usuario sepa si lo que dice se está capturando correctamente.
\end{itemize}

\subsection{CONCLUSIÓN DE LA VALIDACIÓN}

Considerando que se cumplen ambos umbrales del estándar mínimo definido —un 57,9\% de tareas con puntaje igual o superior a 3, por sobre el 55\% requerido, y un promedio global de 2,84, por sobre el mínimo de 2,75—, se declara la propuesta de solución VALIDADA.

Es necesario matizar esta conclusión: los resultados no son contundentes, sino que se ubican apenas por sobre el mínimo exigido. Los modos de falla identificados en la sección anterior y las sugerencias reportadas por los propios usuarios se abordan como trabajo futuro.

\newpage
\secnumbersection{CONCLUSIONES}

\subsection{SÍNTESIS DEL TRABAJO REALIZADO}

A partir de la problemática de accesibilidad web para personas con discapacidad visual, validada empíricamente mediante la encuesta propia del Capítulo 1, se planteó la idea de un relator de pantalla: un asistente basado en un modelo de lenguaje visual (VLM) que no existía previamente bajo esta forma, diseñado para coexistir con el lector de pantalla del usuario en lugar de reemplazarlo -entregando la respuesta a través del mismo mecanismo de accesibilidad del sistema operativo que este ya utiliza-, en línea con el principio de uso equitativo del diseño universal.

Esa idea se implementó de punta a punta mediante una extensión de navegador, un servidor de gestión y un servidor de inferencia propio (\textit{albert}), con dos modalidades de interacción (texto y voz). Para seleccionar el modelo VLM del sistema se realizó un proceso de \textit{benchmark} comparativo entre 5 candidatos, evaluado mediante un juez LLM (gemma3-27b) diseñado y calibrado especialmente para esta tarea, validado contra un criterio humano mediante el esquema de \textit{LLM-as-a-judge-evaluation} a lo largo de 23 versiones de \textit{prompt}. Este juez se estancó en un 72,5\% de \textit{exact match} sin acercarse al 100\%, techo que se interpreta como el máximo potencial de esta implementación particular; se detectó además que parte de su mejor resultado histórico dependía de un sesgo de contexto en los ejemplos de su \textit{prompt}, lo que obligó a una segunda ronda de calibración.

Finalmente, se validó el sistema mediante un conjunto de métricas técnicas (exactitud de STT, latencia del servidor de gestión, latencia de captura, pertinencia de las respuestas) y mediante entrevistas con 3 personas ciegas que usaron el sistema en tareas reales sobre sitios web. Las entrevistas mostraron una impresión general positiva -la mayoría de las tareas intentadas se completaron o al menos se lograron parcialmente-, con buen desempeño en tareas descriptivas o de extracción de información, aunque con un desempeño más débil en tareas procedimentales o de navegación y en portales que requieren autenticación. Entre los modos de falla identificados se cuentan un sesgo de perspectiva visual en las respuestas del modelo, alucinaciones cuando el elemento consultado no estaba visible en la captura, descripciones superficiales en sitios con estructuras más complejas, y fallas repetidas del mecanismo de \textit{scroll \& stitch} al capturar ventanas emergentes.

Algunos aspectos, además, quedaron inconclusos al cierre de este documento. La validación con usuarios se realizó con solo 3 participantes, por lo que sus resultados se reportan como casos ilustrativos y no como evidencia estadísticamente representativa. Algunas sugerencias reportadas por los propios participantes -como el contexto de conversación- se plantean como trabajo futuro en la siguiente sección.

\subsection{TRABAJO FUTURO}

A partir de lo anterior, se identifican las siguientes líneas de trabajo futuro:

\begin{itemize}
    \item Ampliar la validación con usuarios a una muestra mayor, que permita un análisis más representativo que los casos ilustrativos reportados en este trabajo.
    \item Corregir la falla del mecanismo de \textit{scroll \& stitch} al capturar ventanas emergentes.
    \item Reducir el sesgo de perspectiva visual y mejorar el desempeño del sistema en tareas procedimentales o de navegación, y en portales que requieren autenticación.
    \item Incorporar contexto de conversación, para permitir preguntas de seguimiento sobre un mismo elemento sin repetir la información ya entregada.
    \item Mejorar el \textit{feedback} de la entrada por voz, para que el usuario sepa si su pregunta se está capturando correctamente.
    \item Desarrollar una versión compatible con dispositivos móviles y aplicaciones, además de los navegadores de escritorio.
    \item Explorar arquitecturas de juez LLM con acceso directo a la imagen evaluada, para reducir la dependencia de un juez \textit{text-only} frente a alucinaciones del VLM que sean difíciles de detectar solo a partir del texto.
\end{itemize}

\newpage
\secnumberlesssection{ANEXOS}

En los Anexos se incluye todo aquel material complementario que no es parte del contenido de los capítulos de la memoria, pero que permiten a un lector contar con un contenido adjunto relacionado con el tema.

\subsection{ENCUESTA SOBRE TECNOLOGÍAS PARA PERSONAS CIEGAS}
\label{anexo:encuesta}

Este anexo detalla el instrumento mencionado en el Capítulo 1 (Definición del problema), aplicado para validar empíricamente la problemática social antes de proponer una solución.

\begin{table}[h]
    \centering
    \caption{\label{table:fichaencuesta} Ficha técnica de la encuesta.} Fuente: Elaboración propia.
    \begin{tabular}{|p{5cm}|p{9cm}|}
        \hline
        Instrumento & ``Encuesta sobre tecnologías para personas ciegas'' \\
        \hline
        Plataforma & Google Forms \\
        \hline
        Validación previa & Pedro Correa Maturana, director del INLAB (Unidad de Inclusión Laboral) de FUNDALURP (Fundación de Lucha contra la Retinitis Pigmentosa) \\
        \hline
        Canal de distribución & Correo electrónico, a personas asistidas por FUNDALURP o vinculadas a ella \\
        \hline
        Período de aplicación & 26 de febrero de 2026 -- 11 de marzo de 2026 \\
        \hline
        Respuestas obtenidas & 9 (N=9) \\
        \hline
        Encuestados que usan lector de pantalla & 7 de 9 (78\%) \\
        \hline
    \end{tabular}
\end{table}

La encuesta se estructuró en cuatro bloques temáticos:

\begin{itemize}
    \item \textbf{Datos generales y de contacto}: marca temporal de la respuesta, correo electrónico opcional (para diferenciar participantes y evitar respuestas duplicadas), e interés en ser contactado a futuro en el contexto del proyecto.
    \item \textbf{Perfil del encuestado}: género; año de nacimiento (usado para estimar el contexto tecnológico según la época); condición que causa la discapacidad visual (la más severa o antigua, si el encuestado padece más de una); forma en que la condición afecta la visión; nivel de discapacidad visual; grado de manejo computacional actual. Si la discapacidad fue adquirida en una etapa avanzada de la vida (y no de nacimiento), se agregan dos preguntas adicionales: el grado de manejo computacional y el tipo de tecnologías que dominaba antes de la pérdida de visión.
    \item \textbf{Uso de lectores de pantalla}: si el encuestado conoce qué es y qué hace un lector de pantalla; si emplea uno para utilizar su computador; cuáles conoce (aunque no los use); cuál usa como lector principal; y qué impacto tuvo en su vida comenzar a usarlo. Las preguntas de este bloque y del siguiente solo se muestran a quienes reportan emplear un lector de pantalla.
    \item \textbf{Dificultades y funcionalidades}: dificultades típicas al utilizar un lector de pantalla en general; dificultades típicas al navegar sitios web específicamente con uno; selección, de una lista cerrada, de las 3 funcionalidades que le parecerían más útiles; una pregunta abierta opcional para sugerir una funcionalidad adicional no considerada en la lista; y una pregunta abierta opcional de cierre para comentarios sobre la encuesta.
\end{itemize}

\subsubsection{Respuestas completas}

A continuación se listan las 9 respuestas completas de la encuesta. Cada participante se identifica con un código (P1 a P9) en lugar de su correo electrónico, dato personal que se omite por completo. El orden de los participantes corresponde al orden cronológico de respuesta.

Dos preguntas del formulario usan categorías cerradas cuyo texto completo se abrevia en la tabla para no repetirlo en cada fila; su definición completa es la siguiente:

\begin{itemize}
    \item \textbf{Nivel de discapacidad visual}: \textit{Moderada} (dificultad para leer textos estándar y para reconocer rostros a corta distancia, incluso con corrección); \textit{Severa} (no logra leer textos impresos convencionales y requiere apoyos específicos para la orientación y el desplazamiento en espacios públicos); \textit{Ceguera total} (no presenta visión funcional y no puede leer textos sin el uso de apoyos como lectores de pantalla u otros sistemas de accesibilidad).
    \item \textbf{Manejo computacional}: \textit{Básico} (utiliza el computador para navegar por la web, enviar correos y acceder a redes sociales); \textit{Intermedio} (utiliza el computador de forma cotidiana y maneja herramientas de productividad como Word y Excel); \textit{Avanzado} (utiliza software especializado para el trabajo y realiza ajustes en configuraciones del sistema o de las aplicaciones); \textit{Experto} (domina el uso del computador y es capaz de enseñar a otras personas cómo utilizarlo).
\end{itemize}

\begin{longtable}{|p{4cm}|p{10.5cm}|}
    \caption{\label{table:respuestascompletas} Respuestas completas de la encuesta, anonimizadas. Fuente: Elaboración propia.} \\
    \hline
    \textbf{Campo} & \textbf{Respuesta} \\
    \hline
    \endfirsthead
    \hline
    \textbf{Campo} & \textbf{Respuesta} \\
    \hline
    \endhead
    \hline
    \endfoot

    \multicolumn{2}{|l|}{\textbf{Participante P1}} \\
    \hline
    Género & Femenino \\
    \hline
    Año de nacimiento & 1974 \\
    \hline
    Condición & Retinitis pigmentaria \\
    \hline
    Cómo afecta la visión & Ceguera total \\
    \hline
    Nivel de discapacidad visual & Ceguera total \\
    \hline
    Manejo computacional actual & Avanzado \\
    \hline
    Origen de la discapacidad & Adquirida (por enfermedad o accidente) \\
    \hline
    Manejo computacional antes de la pérdida & Experto \\
    \hline
    Tecnologías dominadas antes de la pérdida & Navegar por sitios web, redes sociales (Facebook, Twitter, Instagram), correo electrónico (Gmail, Outlook), tiendas en línea (Falabella, AliExpress), banca en línea, videoconferencias (Zoom, Google Meet), productividad (Word, Excel), manejo de carpetas y archivos del sistema, programas de diseño gráfico, edición audiovisual o arquitectura (Adobe, AutoDesk) \\
    \hline
    ¿Conoce qué es un lector de pantalla? & Sí \\
    \hline
    ¿Emplea un lector de pantalla? & Sí \\
    \hline
    Lectores de pantalla que conoce & JAWS, NVDA, TalkBack, Narrador \\
    \hline
    Lector de pantalla principal & NVDA \\
    \hline
    Impacto de comenzar a usarlo & (sin respuesta) \\
    \hline
    Dificultades al utilizar un lector de pantalla & Incapacidad de interpretar elementos estrictamente visuales como imágenes; incapacidad de leer notificaciones o \textit{pop-ups} independientes de la estructura del sitio o ubicados en otra sección de la página que no está seleccionada en ese momento \\
    \hline
    Dificultades al navegar sitios web & Recuadros de texto a rellenar sin etiqueta o que no informan qué información se debe ingresar; obligación de usar el ratón para interactuar con el sitio; pruebas de verificación humana estrictamente visuales \\
    \hline
    Funcionalidades más útiles seleccionadas & Describir e interpretar imágenes, fotografías y gráficos; solicitar por voz la navegación directa a una sección del sitio (ej. ``Llévame al precio''); lectura de notificaciones o etiquetas dinámicas \\
    \hline
    Funcionalidad adicional sugerida & (sin respuesta) \\
    \hline
    Comentario final & (sin respuesta) \\
    \hline

    \multicolumn{2}{|l|}{\textbf{Participante P2}} \\
    \hline
    Género & Femenino \\
    \hline
    Año de nacimiento & 1975 \\
    \hline
    Condición & Retinitis pigmentaria \\
    \hline
    Cómo afecta la visión & Visión de túnel o periférica reducida; pérdida de visión central; sensibilidad a la luz \\
    \hline
    Nivel de discapacidad visual & Severa \\
    \hline
    Manejo computacional actual & Intermedio \\
    \hline
    Origen de la discapacidad & Adquirida (por enfermedad o accidente) \\
    \hline
    Manejo computacional antes de la pérdida & Intermedio \\
    \hline
    Tecnologías dominadas antes de la pérdida & Navegar por sitios web, redes sociales (Facebook, Twitter, Instagram), correo electrónico (Gmail, Outlook), tiendas en línea (Falabella, AliExpress), banca en línea, videoconferencias (Zoom, Google Meet), productividad (Word, Excel), manejo de carpetas y archivos del sistema \\
    \hline
    ¿Conoce qué es un lector de pantalla? & Sí \\
    \hline
    ¿Emplea un lector de pantalla? & No \\
    \hline

    \multicolumn{2}{|l|}{\textbf{Participante P3}} \\
    \hline
    Género & Masculino \\
    \hline
    Año de nacimiento & 1965 \\
    \hline
    Condición & Degeneración Macular Asociada a la Edad (DMAE) \\
    \hline
    Cómo afecta la visión & Visión borrosa o nublada en general; zonas oscuras o manchas flotantes; sensibilidad a la luz \\
    \hline
    Nivel de discapacidad visual & Severa \\
    \hline
    Manejo computacional actual & Básico \\
    \hline
    Origen de la discapacidad & Adquirida (por enfermedad o accidente) \\
    \hline
    Manejo computacional antes de la pérdida & Intermedio \\
    \hline
    Tecnologías dominadas antes de la pérdida & Navegar por sitios web, redes sociales (Facebook, Twitter, Instagram), correo electrónico (Gmail, Outlook), tiendas en línea (Falabella, AliExpress), banca en línea, manejo de carpetas y archivos del sistema, programas de diseño gráfico, edición audiovisual o arquitectura (Adobe, AutoDesk) \\
    \hline
    ¿Conoce qué es un lector de pantalla? & Sí \\
    \hline
    ¿Emplea un lector de pantalla? & Sí \\
    \hline
    Lectores de pantalla que conoce & TalkBack, Narrador \\
    \hline
    Lector de pantalla principal & VoiceOver \\
    \hline
    Impacto de comenzar a usarlo & ``Fue como volver a ser más autónomo'' \\
    \hline
    Dificultades al utilizar un lector de pantalla & Incapacidad de interpretar elementos estrictamente visuales como imágenes; orden de lectura ilógico que obliga a pasar por información irrelevante antes de llegar a la clave; incapacidad de leer notificaciones o \textit{pop-ups} \\
    \hline
    Dificultades al navegar sitios web & Recuadros de texto sin etiqueta; el sistema entrega información valiosa mediante imágenes sin texto alternativo; pruebas de verificación humana estrictamente visuales \\
    \hline
    Funcionalidades más útiles seleccionadas & Obtener un resumen del contenido principal del sitio; poder hacer preguntas sobre el sitio (ej. ``¿Qué puedo hacer en esta página?''); solicitar por voz la navegación directa a una sección del sitio \\
    \hline
    Funcionalidad adicional sugerida & (sin respuesta) \\
    \hline
    Comentario final & (sin respuesta) \\
    \hline

    \multicolumn{2}{|l|}{\textbf{Participante P4}} \\
    \hline
    Género & Masculino \\
    \hline
    Año de nacimiento & 1972 \\
    \hline
    Condición & Nanoftalmo \\
    \hline
    Cómo afecta la visión & Visión borrosa o nublada en general; baja visión, imposibilidad de identificar figuras, letras o detalles \\
    \hline
    Nivel de discapacidad visual & Severa \\
    \hline
    Manejo computacional actual & Básico \\
    \hline
    Origen de la discapacidad & De nacimiento \\
    \hline
    ¿Conoce qué es un lector de pantalla? & No \\
    \hline
    ¿Emplea un lector de pantalla? & No \\
    \hline

    \multicolumn{2}{|l|}{\textbf{Participante P5}} \\
    \hline
    Género & Masculino \\
    \hline
    Año de nacimiento & 1973 \\
    \hline
    Condición & Opacidad de córnea (síndrome de Peters) \\
    \hline
    Cómo afecta la visión & Visión borrosa o nublada en general; distingue colores y bultos, aunque los colores muy oscuros los percibe como negro (ej. distingue semáforos de noche o con clima nublado) \\
    \hline
    Nivel de discapacidad visual & Severa \\
    \hline
    Manejo computacional actual & Básico \\
    \hline
    Origen de la discapacidad & De nacimiento \\
    \hline
    ¿Conoce qué es un lector de pantalla? & Sí \\
    \hline
    ¿Emplea un lector de pantalla? & Sí \\
    \hline
    Lectores de pantalla que conoce & JAWS, NVDA, TalkBack \\
    \hline
    Lector de pantalla principal & NVDA \\
    \hline
    Impacto de comenzar a usarlo & ``Me siento en igualdad de condiciones ante las personas sin discapacidad visual al usar un pc como ellos'' \\
    \hline
    Dificultades al utilizar un lector de pantalla & Incapacidad de interpretar elementos estrictamente visuales como imágenes; orden de lectura ilógico; incapacidad de leer notificaciones o \textit{pop-ups}; nula reacción ante cambios en la presentación del sitio (secciones que aparecen y desaparecen); respuesta libre: ``los gráficos y cuando no especifica los botones cualquieras éstos sean'' \\
    \hline
    Dificultades al navegar sitios web & El sistema entrega información valiosa mediante imágenes sin texto alternativo; instrucciones estrictamente visuales (ej. ``seleccione la opción verde''); la página cambia su estructura sin avisar; obligación de usar el ratón; pruebas de verificación humana estrictamente visuales; respuesta libre: ``los gráficos y botones sin descripción'' \\
    \hline
    Funcionalidades más útiles seleccionadas & Describir e interpretar imágenes, fotografías y gráficos; solicitar por voz la navegación directa a una sección del sitio; lectura de notificaciones o etiquetas dinámicas \\
    \hline
    Funcionalidad adicional sugerida & ``El significado de los botones y gráficos que muestran los símbolos e imágenes pero para el lector no están descritos en forma escrita, como para que nos diga: botón ingresar, la imagen representa tal o cual figura... cuestiones así por el estilo'' \\
    \hline
    Comentario final & ``Descripción de gestos, textos, gráficos e imágenes'' \\
    \hline

    \multicolumn{2}{|l|}{\textbf{Participante P6}} \\
    \hline
    Género & Masculino \\
    \hline
    Año de nacimiento & 2008 \\
    \hline
    Condición & Neuropatía óptica hereditaria de Leber \\
    \hline
    Cómo afecta la visión & Pérdida de visión central; escasa visión periférica y cambio en la percepción de colores \\
    \hline
    Nivel de discapacidad visual & Severa \\
    \hline
    Manejo computacional actual & Básico \\
    \hline
    Origen de la discapacidad & Adquirida (por enfermedad o accidente) \\
    \hline
    Manejo computacional antes de la pérdida & Básico \\
    \hline
    Tecnologías dominadas antes de la pérdida & Navegar por sitios web, videoconferencias (Zoom, Google Meet) \\
    \hline
    ¿Conoce qué es un lector de pantalla? & No \\
    \hline
    ¿Emplea un lector de pantalla? & Sí \\
    \hline
    Lectores de pantalla que conoce & NVDA, TalkBack \\
    \hline
    Lector de pantalla principal & TalkBack \\
    \hline
    Impacto de comenzar a usarlo & ``Me ayudó a adaptarme a esta nueva normalidad de mejor forma'' \\
    \hline
    Dificultades al utilizar un lector de pantalla & Orden de lectura ilógico que obliga a pasar por información irrelevante antes de llegar a la clave \\
    \hline
    Dificultades al navegar sitios web & La página cambia su estructura u organización sin avisar \\
    \hline
    Funcionalidades más útiles seleccionadas & Obtener un resumen del contenido principal del sitio; describir e interpretar imágenes, fotografías y gráficos; solicitar por voz la navegación directa a una sección del sitio \\
    \hline
    Funcionalidad adicional sugerida & (sin respuesta) \\
    \hline
    Comentario final & (sin respuesta) \\
    \hline

    \multicolumn{2}{|l|}{\textbf{Participante P7}} \\
    \hline
    Género & Otro \\
    \hline
    Año de nacimiento & 1982 \\
    \hline
    Condición & Glaucoma \\
    \hline
    Cómo afecta la visión & (respuesta del encuestado poco clara en el formulario original: ``Lucero visión'') \\
    \hline
    Nivel de discapacidad visual & Ceguera total \\
    \hline
    Manejo computacional actual & Intermedio \\
    \hline
    Origen de la discapacidad & De nacimiento \\
    \hline
    ¿Conoce qué es un lector de pantalla? & Sí \\
    \hline
    ¿Emplea un lector de pantalla? & Sí \\
    \hline
    Lectores de pantalla que conoce & JAWS, NVDA, VoiceOver, TalkBack, Narrador, Orca \\
    \hline
    Lector de pantalla principal & VoiceOver \\
    \hline
    Impacto de comenzar a usarlo & ``Es mi forma y manera de comunicarme con otros y ser productivo laboralmente'' \\
    \hline
    Dificultades al utilizar un lector de pantalla & Incapacidad de interpretar elementos estrictamente visuales como imágenes; orden de lectura ilógico; incapacidad de leer notificaciones o \textit{pop-ups}; nula reacción ante cambios en la presentación del sitio \\
    \hline
    Dificultades al navegar sitios web & Recuadros de texto sin etiqueta; imágenes sin texto alternativo; navegar textos extensos sin estructura; instrucciones estrictamente visuales; la página cambia su estructura sin avisar; obligación de usar el ratón; pruebas de verificación humana estrictamente visuales \\
    \hline
    Funcionalidades más útiles seleccionadas & Obtener un resumen del contenido principal del sitio; poder hacer preguntas sobre el sitio; describir e interpretar imágenes, fotografías y gráficos \\
    \hline
    Funcionalidad adicional sugerida & (sin respuesta) \\
    \hline
    Comentario final & (sin respuesta) \\
    \hline

    \multicolumn{2}{|l|}{\textbf{Participante P8}} \\
    \hline
    Género & Femenino \\
    \hline
    Año de nacimiento & 1979 \\
    \hline
    Condición & Glaucoma \\
    \hline
    Cómo afecta la visión & Ceguera total \\
    \hline
    Nivel de discapacidad visual & Ceguera total \\
    \hline
    Manejo computacional actual & Básico \\
    \hline
    Origen de la discapacidad & Adquirida (por enfermedad o accidente) \\
    \hline
    Manejo computacional antes de la pérdida & Intermedio \\
    \hline
    Tecnologías dominadas antes de la pérdida & Navegar por sitios web \\
    \hline
    ¿Conoce qué es un lector de pantalla? & Sí \\
    \hline
    ¿Emplea un lector de pantalla? & Sí \\
    \hline
    Lectores de pantalla que conoce & JAWS, NVDA, VoiceOver, TalkBack \\
    \hline
    Lector de pantalla principal & NVDA \\
    \hline
    Impacto de comenzar a usarlo & (sin respuesta) \\
    \hline
    Dificultades al utilizar un lector de pantalla & Incapacidad de interpretar elementos estrictamente visuales como imágenes; orden de lectura ilógico; incapacidad de leer notificaciones o \textit{pop-ups} \\
    \hline
    Dificultades al navegar sitios web & Recuadros de texto sin etiqueta; imágenes sin texto alternativo; navegar textos extensos sin estructura; instrucciones estrictamente visuales; pruebas de verificación humana estrictamente visuales \\
    \hline
    Funcionalidades más útiles seleccionadas & Obtener un resumen del contenido principal del sitio; describir e interpretar imágenes, fotografías y gráficos; solicitar por voz la navegación directa a una sección del sitio \\
    \hline
    Funcionalidad adicional sugerida & (sin respuesta) \\
    \hline
    Comentario final & (sin respuesta) \\
    \hline

    \multicolumn{2}{|l|}{\textbf{Participante P9}} \\
    \hline
    Género & Femenino \\
    \hline
    Año de nacimiento & 1967 \\
    \hline
    Condición & Retinitis pigmentaria \\
    \hline
    Cómo afecta la visión & Visión de túnel o periférica reducida; zonas oscuras o manchas flotantes \\
    \hline
    Nivel de discapacidad visual & Moderada \\
    \hline
    Manejo computacional actual & Básico \\
    \hline
    Origen de la discapacidad & Adquirida (por enfermedad o accidente) \\
    \hline
    Manejo computacional antes de la pérdida & Básico \\
    \hline
    Tecnologías dominadas antes de la pérdida & Productividad (Word, Excel) \\
    \hline
    ¿Conoce qué es un lector de pantalla? & Sí \\
    \hline
    ¿Emplea un lector de pantalla? & Sí \\
    \hline
    Lectores de pantalla que conoce & NVDA, TalkBack \\
    \hline
    Lector de pantalla principal & NVDA \\
    \hline
    Impacto de comenzar a usarlo & ``Me emocioné mucho porque, cuando escribo, el lector me va leyendo y así estoy segura de lo que voy escribiendo'' \\
    \hline
    Dificultades al utilizar un lector de pantalla & Respuesta libre: ``mi dificultad es que el lector no deja de hablar'' \\
    \hline
    Dificultades al navegar sitios web & Navegar textos extensos sin estructura; instrucciones estrictamente visuales \\
    \hline
    Funcionalidades más útiles seleccionadas & Poder hacer preguntas sobre el sitio; describir e interpretar imágenes, fotografías y gráficos; solicitar por voz la navegación directa a una sección del sitio \\
    \hline
    Funcionalidad adicional sugerida & (sin respuesta) \\
    \hline
    Comentario final & ``Me parece una iniciativa muy buena que se sigan buscando mejoras para el lector de pantalla'' \\
    \hline
\end{longtable}

\subsection{CALIBRACIÓN DEL JUEZ LLM: HISTORIAL DE VERSIONES DEL \textit{PROMPT}}
\label{anexo:promptsjuez}

Este anexo detalla el proceso mencionado en el Capítulo 3 (sección \ref{sec:benchmark}, ``Selección del juez''): la calibración del \textit{prompt} del juez LLM contra el conjunto de 40 respuestas evaluadas por un humano, a lo largo de 23 versiones (v1 a v23). Para el candidato gemma3-27b se incluye además el desglose de aciertos por nivel de la rúbrica (L0/L1/L2/L3, sobre un total de 10 casos por nivel).

\begin{longtable}{|p{1.5cm}|p{5cm}|p{2.8cm}|p{2.8cm}|}
    \caption{\label{table:historialprompts} Resultados de \textit{exact match} por versión de \textit{prompt} y modelo juez. Fuente: Elaboración propia.} \\
    \hline
    \textbf{Versión} & \textbf{gemma3-27b (L0/L1/L2/L3)} & \textbf{gemma-3-12b} & \textbf{MPrometheus-14B} \\
    \hline
    \endfirsthead
    \hline
    \textbf{Versión} & \textbf{gemma3-27b (L0/L1/L2/L3)} & \textbf{gemma-3-12b} & \textbf{MPrometheus-14B} \\
    \hline
    \endhead
    \hline
    \endfoot
    v1 & sin métricas formales & -- & -- \\
    \hline
    v2 & sin métricas formales & -- & -- \\
    \hline
    v3 & sin métricas formales & -- & -- \\
    \hline
    v4 & sin métricas formales & -- & -- \\
    \hline
    v5 & 70,0\% (7/3/9/9) & 60,0\% & 45,0\% \\
    \hline
    v6 & 70,0\% (4/6/8/10) & 57,5\% & 52,5\% \\
    \hline
    v7 & 67,5\% (4/6/7/10) & 52,5\% & 45,0\% \\
    \hline
    v8 & 62,5\% (3/5/7/10) & 55,0\% & 50,0\% \\
    \hline
    v9 & \textbf{72,5\% (5/6/8/10)} & 60,0\% & 52,5\% \\
    \hline
    v10 & 72,5\% (4/6/9/10) & 55,0\% & 42,5\% \\
    \hline
    v11 & 70,0\% (5/5/8/10) & 62,5\% & 57,5\% \\
    \hline
    v12 & 70,0\% (4/5/9/10) & 57,5\% & 47,5\% \\
    \hline
    v13 & 70,0\% (5/5/8/10) & -- & -- \\
    \hline
    v14 & 67,5\% (3/6/8/10) & 50,0\% & 45,0\% \\
    \hline
    v15 & 42,5\% (5/8/2/2) & 37,5\% & 32,5\% \\
    \hline
    v16 & 70,0\% (6/4/8/10) & -- & -- \\
    \hline
    v17 & 70,0\% (6/4/8/10) & -- & -- \\
    \hline
    v18 & 70,0\% (6/4/8/10) & -- & -- \\
    \hline
    v19 & 67,5\% (6/3/8/10) & -- & -- \\
    \hline
    v20 & \textbf{72,5\% (6/5/8/10)} & -- & -- \\
    \hline
    v21 & 70,0\% (6/6/6/10) & -- & -- \\
    \hline
    v22 & 67,5\% (6/4/7/10) & -- & -- \\
    \hline
    v23 & 70,0\% (6/6/6/10) & -- & -- \\
    \hline
\end{longtable}

\subsection{DETALLE DE LAS SESIONES DE ENTREVISTAS}
\label{anexo:sesiones}

Este anexo detalla el proceso mencionado en el Capítulo 4 (sección de entrevistas con usuarios): las 4 sesiones realizadas sobre los 3 participantes, con su modalidad, duración y los sitios evaluados en cada una.

\begin{table}[h]
    \centering
    \caption{\label{table:detallesesiones} Detalle de las sesiones de entrevistas.} Fuente: Elaboración propia.
    \begin{tabular}{|p{1.15cm}|p{1.8cm}|p{1.9cm}|p{1.6cm}|p{4.75cm}|p{1.4cm}|}
        \hline
        \textbf{Sesión} & \textbf{Usuario} & \textbf{Modalidad} & \textbf{Duración} & \textbf{Sitios evaluados} & \textbf{N° tareas} \\
        \hline
        0 & Usuario 0 & Texto & 30 min & USM (educacional), Falabella (\textit{e-commerce}) & 4 \\
        \hline
        1 & Usuario 1 & Texto & 1 hora & YouTube (redes sociales), portal docente de una institución educativa (educacional) & 9 \\
        \hline
        2 & Usuario 2 & Texto & 15 min & Portal docente y portal de estudiante de instituciones educativas (educacional), SII (estatal) & 5 \\
        \hline
        3 & Usuario 2 & Voz & 15 min & SII (estatal) & 1 \\
        \hline
    \end{tabular}
\end{table}

\subsection{REGISTRO DE LA PRUEBA DE RECONOCIMIENTO DE VOZ (STT)}
\label{anexo:stt}

Este anexo detalla el proceso mencionado en el Capítulo 4 (sección de métricas técnicas, exactitud del reconocimiento de voz): las 10 frases leídas en voz alta por cada una de las 5 personas que participaron en la prueba, la transcripción literal devuelta por el motor de STT de la extensión, y la cantidad de palabras erróneas de cada una respecto del total de palabras de la frase original.

\begin{longtable}{|p{0.5cm}|p{5.7cm}|p{5.7cm}|p{1.5cm}|}
    \caption{\label{table:transcripcionesstt} Registro completo de la prueba de reconocimiento de voz: frase original, transcripción obtenida y palabras erróneas por frase, para cada una de las 5 personas.} \\
    \hline
    \textbf{\#} & \textbf{Frase original} & \textbf{Transcripción obtenida} & \textbf{Palabras erróneas} \\
    \hline
    \endfirsthead
    \hline
    \textbf{\#} & \textbf{Frase original} & \textbf{Transcripción obtenida} & \textbf{Palabras erróneas} \\
    \hline
    \endhead
    \hline
    \endfoot

    \multicolumn{4}{|l|}{\textbf{Persona 1}} \\
    \hline
    1 & ¿Dónde puedo encontrar los canales de contacto oficial de la institución? & Dónde puedo encontrar los canales de contacto oficial de la institución & 0 / 11 \\
    \hline
    2 & Indícame, de haber, dónde debo acceder para ingresar a mi cuenta. & Indícame a ver dónde debo acceder para ingresar a mi cuenta & 2 / 11 \\
    \hline
    3 & ¿En qué parte exacta de la pantalla debo ubicarme para revisar el carrito de compras? & En qué parte exacta la pantalla debe ubicarme para revisar el cargo de compras & 3 / 15 \\
    \hline
    4 & ¿Qué forma tienen los elementos interactuables ubicados debajo de cada publicación, y qué propósito tienen? & Qué forma tienen los elementos intelectual de ubicados debajo de cada publicación Y qué propósito tiene & 3 / 15 \\
    \hline
    5 & ¿Cuál es la temática de este sitio web, y qué tipo de contenido se visualiza en él? & Cuál es la temática de este sitio web Y qué tipo de contenido se visualiza en él & 0 / 17 \\
    \hline
    6 & Basado en la forma estructural del sitio, ¿cómo se interactúa físicamente con la página para seguir consumiendo contenido? & basado en la forma estructural del sitio Cómo se interactúa físicamente con la página para seguir consumiendo contenido & 0 / 18 \\
    \hline
    7 & ¿Dónde se encuentra el menú principal de navegación y qué opciones o portales clave destaca para estudiantes o postulantes? & Dónde se encuentra el menú principal de navegación y Qué opciones o portales clave destaca para sus estudiantes de postulante & 3 / 19 \\
    \hline
    8 & ¿Cuál es el propósito principal de esta página web y qué tipo de información o servicios públicos se ofrecen en ella? & Cuál es el propósito principal de esta página web Y qué tipo de información o servicios públicos se ofrecen en ella & 0 / 21 \\
    \hline
    9 & ¿De qué trata esta página web y cuáles son los bloques o secciones que componen su estructura principal de arriba hacia abajo? & De qué trata esta página web y cuáles son los bloqueos secciones que componen su estructura principal de arriba hacia abajo & 2 / 22 \\
    \hline
    10 & ¿Dónde se ubican espacialmente las herramientas de accesibilidad del sitio (si las hay) y qué funciones de ajuste visual o lectura ofrecen? & Dónde se ubican espacialmente la herramienta de accesibilidad del sitio si las hay y qué funciones de ajuste visual la lectura ofrecen & 3 / 22 \\
    \hline

    \multicolumn{4}{|l|}{\textbf{Persona 2}} \\
    \hline
    1 & ¿Dónde puedo encontrar los canales de contacto oficial de la institución? & Dónde puedo encontrar los canales de contacto oficial de la institución & 0 / 11 \\
    \hline
    2 & Indícame, de haber, dónde debo acceder para ingresar a mi cuenta. & Indícame a ver dónde puedo acceder para ingresar a mi cuenta & 3 / 11 \\
    \hline
    3 & ¿En qué parte exacta de la pantalla debo ubicarme para revisar el carrito de compras? & En qué parte exacta de la pantalla debo ubicarme para revisar el carrito de compras & 0 / 15 \\
    \hline
    4 & ¿Qué forma tienen los elementos interactuables ubicados debajo de cada publicación, y qué propósito tienen? & Qué forma tiene los elementos interactuables ubicados debajo de cada publicación Y qué propósito tienen & 1 / 15 \\
    \hline
    5 & ¿Cuál es la temática de este sitio web, y qué tipo de contenido se visualiza en él? & Cuál es la temática de este sitio web Y qué tipo de contenido se visualiza en él & 0 / 17 \\
    \hline
    6 & Basado en la forma estructural del sitio, ¿cómo se interactúa físicamente con la página para seguir consumiendo contenido? & basado en la fórmula estructural del sirio Cómo se interactúa físicamente con la página para seguir consumiendo contenido & 2 / 18 \\
    \hline
    7 & ¿Dónde se encuentra el menú principal de navegación y qué opciones o portales clave destaca para estudiantes o postulantes? & Dónde se encuentra el menú principal de navegación y Qué opciones o portales clave destaca para estudiantes o postulantes & 0 / 19 \\
    \hline
    8 & ¿Cuál es el propósito principal de esta página web y qué tipo de información o servicios públicos se ofrecen en ella? & Cuál es el propósito principal de esta página web Y qué tipo de información o servicios públicos se ofrecen en ella & 0 / 21 \\
    \hline
    9 & ¿De qué trata esta página web y cuáles son los bloques o secciones que componen su estructura principal de arriba hacia abajo? & De qué trata esta página web y cuáles son los bloques o secciones que componen su estructura principal de arriba hacia abajo & 0 / 22 \\
    \hline
    10 & ¿Dónde se ubican espacialmente las herramientas de accesibilidad del sitio (si las hay) y qué funciones de ajuste visual o lectura ofrecen? & Dónde se ubican especialmente las herramientas de accesibilidad del sitio y Qué funciones de ajuste visual o lectura ofrecen & 4 / 22 \\
    \hline

    \multicolumn{4}{|l|}{\textbf{Persona 3}} \\
    \hline
    1 & ¿Dónde puedo encontrar los canales de contacto oficial de la institución? & Dónde puedo encontrar los canales de contacto oficial de la institución & 0 / 11 \\
    \hline
    2 & Indícame, de haber, dónde debo acceder para ingresar a mi cuenta. & Indícame de haber dónde debo acceder para ingresar a mi cuenta & 0 / 11 \\
    \hline
    3 & ¿En qué parte exacta de la pantalla debo ubicarme para revisar el carrito de compras? & En qué parte exacta de la pantalla debo ubicarme para revisar el carrito de compras & 0 / 15 \\
    \hline
    4 & ¿Qué forma tienen los elementos interactuables ubicados debajo de cada publicación, y qué propósito tienen? & Qué forma tienen los elementos interactuables ubicados debajo de cada publicación Y qué propósito tienen & 0 / 15 \\
    \hline
    5 & ¿Cuál es la temática de este sitio web, y qué tipo de contenido se visualiza en él? & Cuál es la temática de este sitio web Y qué tipo de contenido se visualiza en él & 0 / 17 \\
    \hline
    6 & Basado en la forma estructural del sitio, ¿cómo se interactúa físicamente con la página para seguir consumiendo contenido? & basado en la forma estructural del sitio Cómo se interactúa físicamente con la página para seguir consumiendo contenido & 0 / 18 \\
    \hline
    7 & ¿Dónde se encuentra el menú principal de navegación y qué opciones o portales clave destaca para estudiantes o postulantes? & Dónde se encuentra el menú principal de navegación y Qué opciones soportales clave destaca para estudiantes o postulantes & 2 / 19 \\
    \hline
    8 & ¿Cuál es el propósito principal de esta página web y qué tipo de información o servicios públicos se ofrecen en ella? & Cuál es el propósito principal de esta página web Y qué tipo de información o servicios públicos se ofrecen en ella & 0 / 21 \\
    \hline
    9 & ¿De qué trata esta página web y cuáles son los bloques o secciones que componen su estructura principal de arriba hacia abajo? & De qué trata esta página web y cuáles son los bloques o secciones que componen su estructura principal de arriba hacia abajo & 0 / 22 \\
    \hline
    10 & ¿Dónde se ubican espacialmente las herramientas de accesibilidad del sitio (si las hay) y qué funciones de ajuste visual o lectura ofrecen? & Dónde se ubican especialmente las herramientas de accesibilidad del sitio si las hay y qué funciones de ajuste visual o lectura ofrecen & 1 / 22 \\
    \hline

    \multicolumn{4}{|l|}{\textbf{Persona 4}} \\
    \hline
    1 & ¿Dónde puedo encontrar los canales de contacto oficial de la institución? & Dónde puedo encontrar los canales de contacto oficial de la institución & 0 / 11 \\
    \hline
    2 & Indícame, de haber, dónde debo acceder para ingresar a mi cuenta. & Indícame de a ver dónde debo acceder para ingresar a mi cuenta & 2 / 11 \\
    \hline
    3 & ¿En qué parte exacta de la pantalla debo ubicarme para revisar el carrito de compras? & En qué parte exacta de la pantalla debo ubicarme para revisar el carrito de compras & 0 / 15 \\
    \hline
    4 & ¿Qué forma tienen los elementos interactuables ubicados debajo de cada publicación, y qué propósito tienen? & Qué forma tienen los elementos interactuables ubicados debajo de cada publicación Y qué propósito tienen & 0 / 15 \\
    \hline
    5 & ¿Cuál es la temática de este sitio web, y qué tipo de contenido se visualiza en él? & Cuál es la temática de este sitio web Y qué tipo de contenido se visualiza en él & 0 / 17 \\
    \hline
    6 & Basado en la forma estructural del sitio, ¿cómo se interactúa físicamente con la página para seguir consumiendo contenido? & basado en la forma estructural del sitio Cómo se interactúa físicamente con la página para seguir consumiendo contenido & 0 / 18 \\
    \hline
    7 & ¿Dónde se encuentra el menú principal de navegación y qué opciones o portales clave destaca para estudiantes o postulantes? & Dónde se encuentra el menú principal de navegación y Qué opciones o portales clave destaca para estudiantes o postulantes & 0 / 19 \\
    \hline
    8 & ¿Cuál es el propósito principal de esta página web y qué tipo de información o servicios públicos se ofrecen en ella? & Cuál es el propósito principal de esta página web Y qué tipo de información o servicios públicos se ofrecen en ella & 0 / 21 \\
    \hline
    9 & ¿De qué trata esta página web y cuáles son los bloques o secciones que componen su estructura principal de arriba hacia abajo? & De qué trata esta página web y cuáles son los bloques o secciones que componen su estructura principal de arriba hacia abajo & 0 / 22 \\
    \hline
    10 & ¿Dónde se ubican espacialmente las herramientas de accesibilidad del sitio (si las hay) y qué funciones de ajuste visual o lectura ofrecen? & Dónde se ubican espacialmente las herramientas de accesibilidad del sitio si las hay y qué funciones de ajuste visual o lectura ofrecen & 0 / 22 \\
    \hline

    \multicolumn{4}{|l|}{\textbf{Persona 5}} \\
    \hline
    1 & ¿Dónde puedo encontrar los canales de contacto oficial de la institución? & Dónde puedo encontrar los canales de contacto oficial de la institución & 0 / 11 \\
    \hline
    2 & Indícame, de haber, dónde debo acceder para ingresar a mi cuenta. & Indícame a ver dónde debo acceder para ingresar a mi cuenta & 2 / 11 \\
    \hline
    3 & ¿En qué parte exacta de la pantalla debo ubicarme para revisar el carrito de compras? & En qué parte salta de la pantalla de ubicarme para revisar el carrito de compras & 2 / 15 \\
    \hline
    4 & ¿Qué forma tienen los elementos interactuables ubicados debajo de cada publicación, y qué propósito tienen? & Qué forma tiene los elementos intelectuales ubicados debajo de cada publicación Y qué propósito tienen & 2 / 15 \\
    \hline
    5 & ¿Cuál es la temática de este sitio web, y qué tipo de contenido se visualiza en él? & Cuál es la temática de este sitio web Y qué tipo de contenido se visualiza en él & 0 / 17 \\
    \hline
    6 & Basado en la forma estructural del sitio, ¿cómo se interactúa físicamente con la página para seguir consumiendo contenido? & basado en la forma estructural del sitio Cómo se interactúa físicamente con la página para seguir consumiendo contenido & 0 / 18 \\
    \hline
    7 & ¿Dónde se encuentra el menú principal de navegación y qué opciones o portales clave destaca para estudiantes o postulantes? & Dónde se encuentra el menú principal de navegación y Qué opciones o portales clave destaca para estudiantes y postulantes & 1 / 19 \\
    \hline
    8 & ¿Cuál es el propósito principal de esta página web y qué tipo de información o servicios públicos se ofrecen en ella? & Cuál es el propósito principal de esta página web Y qué tipo de información o servicios públicos se ofrecen en ella & 0 / 21 \\
    \hline
    9 & ¿De qué trata esta página web y cuáles son los bloques o secciones que componen su estructura principal de arriba hacia abajo? & De qué trata esta página web y cuáles son los bloques o secciones que componen su estructura principal de arriba hacia abajo & 0 / 22 \\
    \hline
    10 & ¿Dónde se ubican espacialmente las herramientas de accesibilidad del sitio (si las hay) y qué funciones de ajuste visual o lectura ofrecen? & Dónde se ubican especialmente las herramientas de accesibilidad del sitio y Qué funciones de ajuste visual o lectura ofrecen & 4 / 22 \\
    \hline
\end{longtable}


\newpage
% Bibliografía estilo APA:
\bibliographystyle{apalike-es}
\bibliography{bibliografia}{}

\end{document}
